% Korean Noether paper 19 — fresh producer translation from ED0005.
% Initial translation source: ED0005 cumulative lines 11243--12397.
% Current scope-rebase authority: ED0011 / NOETH-DE-AUTH-v047-20260808.
% Producer state only; independent checking and certification remain external.

\begin{center}
{\Large\bfseries 환영역의 아이디얼론.}\par
\vspace{0.7em}
지은이\par
\vspace{0.25em}
괴팅겐의 에미 뇌터.\par
\vspace{0.9em}
\rule{4em}{0.4pt}\par
\vspace{0.9em}
{\large\bfseries 목차.}
\end{center}
\noindent 서론.\par
\noindent\hangindent=3.0em\hangafter=1 \S~1. \textbf{환영역, 아이디얼, 유한조건.}\par
\noindent\hangindent=3.0em\hangafter=1 \S~2. \textbf{한 아이디얼을 유한개의 기약아이디얼의 최소공배수로 나타내기.}\par
\noindent\hangindent=3.0em\hangafter=1 \S~3. \textbf{기약아이디얼로의 서로 다른 두 분해에서 성분수의 일치.}\par
\noindent\hangindent=3.0em\hangafter=1 \S~4. \textbf{으뜸아이디얼. 기약아이디얼로의 서로 다른 두 분해에서 대응 소아이디얼의 유일성.}\par
\noindent\hangindent=3.0em\hangafter=1 \S~5. \textbf{한 아이디얼을 극대 으뜸아이디얼들의 최소공배수로 나타내기. 대응 소아이디얼의 유일성.}\par
\noindent\hangindent=3.0em\hangafter=1 \S~6. \textbf{한 아이디얼을 상대소-기약 아이디얼들의 최소공배수로 유일하게 나타내기.}\par
\noindent\hangindent=3.0em\hangafter=1 \S~7. \textbf{고립아이디얼의 유일성.}\par
\noindent\hangindent=3.0em\hangafter=1 \S~8. \textbf{한 아이디얼을 서로소-기약 아이디얼들의 곱으로 유일하게 나타내기.}\par
\noindent\hangindent=3.0em\hangafter=1 \S~9. \textbf{모듈로의 고찰 확장. 기약모듈로의 분해에서 성분수의 일치.}\par
\noindent\hangindent=3.0em\hangafter=1 \S~10. \textbf{다항식영역의 특수한 경우.}\par
\noindent\hangindent=3.0em\hangafter=1 \S~11. \textbf{수론과 미분식 이론의 예.}\par
\noindent\hangindent=3.0em\hangafter=1 \S~12. \textbf{초등인자론의 예.}\par

\begin{center}\textbf{서론.}\end{center}
이 논문의 내용은 \emph{유리정수의 분해정리, 또는 대수적 수체에서 아이디얼의 분해정리를 임의의 정역, 더 일반적으로는 환영역의 아이디얼로 옮기는 것}이다. 이 옮김을 이해하기 위해 먼저 유리정수의 분해정리를 통상적인 정식화와 조금 다른 꼴로 적어 두자.

식
\[
 a=p_1^{\varrho_1}p_2^{\varrho_2}\cdots p_\sigma^{\varrho_\sigma}=q_1q_2\cdots q_\sigma
\]
에서 소수의 거듭제곱 $q_i$들을 분해의 성분으로 보면, 이 성분들은 다음과 같은 특징을 갖는다.

1. 이들은 \emph{쌍마다 서로소}이다. 그러나 어느 $q_i$도 쌍마다 서로소인 수들의 곱으로 나타낼 수 없으므로, 이 뜻에서 기약성도 성립한다. 쌍마다 서로소라는 사실에서 곱 $q_1\cdots q_\sigma$가 최소공배수 $[q_1\cdots q_\sigma]$와 같다는 것도 따른다.

2. 임의의 두 성분 $q_i,q_k$는 \emph{상대소}이다. 곧 $bq_i$가 $q_k$로 나누어지면 $b$가 $q_k$로 나누어진다. 이 뜻에서도 기약성이 성립한다.

3. 각 $q$는 \emph{으뜸}이다. 곧 곱 $b\cdot c$가 $q$로 나누어지지만 $b$는 나누어지지 않으면, $c$의 어떤 거듭제곱\footnote{이 거듭제곱을 언제나 첫째 거듭제곱으로 잡을 수 있으면, 잘 알려진 대로 소수를 얻는다.}이 $q$로 나누어진다. 더 나아가 서로 다른 두 $q$의 곱은 더 이상 으뜸이 아니므로, 이 표현은 \emph{극대 으뜸성분들}에 의한 표현이다. 극대 으뜸성분들로의 분해에 관해서도 $q$들은 기약이다.

4. 각 $q$는 두 진약수의 최소공배수로 나타낼 수 없다는 뜻에서 \emph{기약}이다.

이 으뜸수 $q$와 소수 $p$의 관계는 다음과 같다. 각 $q$마다 $q$의 약수이면서 어떤 거듭제곱이 $q$로 나누어지는 $p$가 하나, 그리고 부호를 제외하면 오직 하나 존재한다. 이것이 대응 소수이다. $p^\varrho$가 그런 거듭제곱 가운데 가장 낮은 것, 곧 $\varrho$가 $q$의 지수이면, 여기서는 특히 $p^\varrho=q$이다. 이제 \emph{유일성 정리}는 다음과 같이 말할 수 있다.

\emph{한 유리정수를 기약인 극대 으뜸성분 $q$들로 나타내는 서로 다른 두 분해에서는 성분수, 대응 소수들(부호를 제외하고), 그리고 지수들이 일치한다. $p^\varrho=q$이므로 $q$들 자체도 부호를 제외하고 일치한다.}

수 대신 그 수에서 유도되는 아이디얼, 곧 $a$로 나누어지는 모든 수를 고찰하면 부호에서 생기는 부정성이 사라진다는 것은 잘 알려져 있다. 그러면 같은 정식화가 (유한) 대수적 수체의 아이디얼을 소아이디얼의 거듭제곱으로 유일하게 분해하는 정리에 그대로 적용된다.

이하에서는 (\S~1) 각 아이디얼이 유한 아이디얼기저를 갖는다는 \emph{유한조건}만을 만족하는 일반 환영역을 바탕으로 삼는다. 이러한 유한조건이 없으면 기약아이디얼과 소아이디얼이 존재하지 않을 수도 있다. 실제로 \emph{모든 대수적 정수}의 영역에는 \emph{소아이디얼로의 분해가 없다}.

성분 $q$의 네 가지 특징에 대응하여 일반적으로는 \emph{서로 구별되는 네 분해가 존재하며, 각 분해는 앞의 분해를 더 세분하여 얻어짐}이 드러난다. 서로소-기약 아이디얼로의 분해는 곱표현이고, 나머지 세 분해는 최소공배수에 의한 축약표현(\S~2)이다. 으뜸아이디얼과 대응 소아이디얼 사이의 \emph{관계}도 유지된다. 기약아이디얼도 으뜸아이디얼이다. 각 으뜸아이디얼 $\ideal Q$에는 $\ideal Q$의 약수이고 그 어떤 거듭제곱이 $\ideal Q$로 나누어지는 대응 소아이디얼 $\ideal P$가 유일하게 정해진다. $\ideal P^\varrho$가 그러한 거듭제곱 가운데 가장 낮은 것, 곧 $\varrho$가 $\ideal Q$의 지수라고 해도 여기서는 $\ideal P^\varrho$와 $\ideal Q$가 일치할 필요가 없다. 이제 \emph{유일성 정리}는 다음과 같이 말한다.

\emph{분해 1과 2는 유일하다. 분해 3 또는 4의 서로 다른 두 표현에서는 성분수와 대응 소아이디얼들이 일치한다}\footnote{아마 지수들도 일치하며, 더 일반적으로 대응 성분들이 동형이라는 사실도 성립할 것이다.}\emph{. 성분들 가운데 나타나는 고립아이디얼(\S~7)은 유일하게 정해진다.}

분해정리를 증명하기 위해 먼저 유한조건에서 데데킨트가 유한 수모듈에 대하여 처음 말한 ``유한사슬 정리''를 이끌어 내고, 그로부터 각 아이디얼을 유한개의 기약아이디얼의 최소공배수로 나타내는 분해 4를 얻는다. 한 성분의 가약성 개념을 바꾸어 쓰면 \emph{기약아이디얼로의 분해 4에 관한 기본 유일성 정리}가 나온다. 이어서 매번 유한개의 성분을 묶어 나머지 분해들로 올라가며, 그 유일성 정리들은 분해 4의 유일성 정리에서 따라온다.

끝으로 (\S~9) 유한개의 기약성분에 의한 표현은 더 약한 가정 아래에서도 성립함을 보인다. 환영역의 가환성은 필요하지 않고, 아이디얼 대신 그 영역에 관한 모듈을 고찰해도 충분하다. 이 더 일반적인 경우에도 서로 다른 두 분해에서 성분수는 일치한다. 반면 소와 으뜸이라는 개념은 가환성과 아이디얼 개념에 매여 있다. 그러나 비가환 영역의 아이디얼에서는 서로소라는 개념이 그대로 남는다.

실제로 서로 구별되는 네 분해가 나타나는 가장 간단한 환영역은 임의의 복소수 계수를 갖는 $n$변수 다항식 전체의 영역이다. 여기서는 각 분해를 대수적 형상의 거동으로 무리적으로 해석할 수 있고, 대응 소아이디얼에 관한 유일성 정리는 대수적 형상을 기약성분들로 유일하게 분해할 수 있다는 소거이론의 기본정리에 대응한다. 더 일반적인 예는 다항식으로 이루어진 모든 유한 정역이다(\S~10). 그러나 모든 짝수의 단순한 영역, 더 일반적으로 한 고정된 수로 나누어지는 모든 수의 영역에서도 이미 일부 분리된 분해의 예를 볼 수 있다(\S~11). 비가환 영역의 아이디얼론의 한 예는 초등인자론(\S~12)이 준다. 거기서는 아이디얼 또는 클래스가 기약아이디얼 또는 기약클래스로 유일하게 분해된다. 이 기약클래스들은 초등인자의 기약성분을 완전히 특징짓고, 보통의 초등인자론이 성립하지 않는 영역에서는 그 등가물로 볼 수도 있을 것이다.

기존 문헌에 관해서는 다음을 말해 둘 필요가 있다. 임의의 복소수 계수 또는 정수 계수를 갖는 다항식영역에서 극대 으뜸아이디얼로의 분해는 라스커가 제시했고, 매콜리가 몇몇 점에서 더 발전시켰다.\footnote{E. Lasker, \emph{Zur Theorie der Moduln und Ideale}. Math. Ann. 60 (1905), S. 20, Satz VII und XIII. -- F. S. Macaulay, \emph{On the Resolution of a given Modular System into Primary Systems including some Properties of Hilbert Numbers}. Math. Ann. 74 (1913), S. 66.} 두 사람은 모두 소거이론, 곧 다항식이 기약다항식들의 곱으로 유일하게 표현된다는 사실을 이용한다. 그러나 대수적 수체의 아이디얼론에서 짐작할 수 있고 이 논문이 보여 주듯, 아이디얼의 분해정리들은 실제로 이 가정과 무관하다. 라스커와 매콜리는 으뜸아이디얼도 소거이론의 개념들을 바탕으로 정의한다.

기약아이디얼로의 분해와 상대소-기약 아이디얼로의 분해는 다항식영역에 대해서도 문헌에서 주목되지 않은 듯하다. 매콜리에게서 고립된 으뜸아이디얼의 유일성에 관한 한 논평만을 찾을 수 있다.

다항식영역에서 서로소-기약 아이디얼로의 분해는 슈마이들러\footnote{W. Schmeidler, \emph{Über Moduln und Gruppen hyperkomplexer Größen}. Math. Zeitschr. 3 (1919), S. 29.}가 유한성 증명에 소거이론을 사용하여 제시했다. 그러나 거기서는 유일성 정리가 아이디얼 자체가 아니라 아이디얼의 클래스에 대해서만 진술된다. 아이디얼 자체에 관한 이 유일성 정리는 비가환 다항식영역의 아이디얼을 다룬 공동논문\footnote{E. Noether -- W. Schmeidler, \emph{Moduln in nichtkommutativen Bereichen, insbesondere aus Differential- und Differenzenausdrücken}. Math. Zeitschr. 8 (1920), S. 1.}에 있다. 그곳에서는 유한 아이디얼기저만을 사용하므로 정리와 방법은 일반 환영역에서도 유지되며, 이 논문에 의해 성분수의 일치라는 점에서 더 날카롭게 만들 수 있다(\S~11). 이 논문의 고찰은 두 앞선 논문의 개념구성을 크게 일반화하고 발전시킨 것이다. 두 논문의 핵심은 최소공배수에 의한 표현에서 잉여류계의 가법적 분해로 넘어가는 데 있다. 여기서는 서술을 간단히 하려고 다시 최소공배수를 사용한다. 그러면 가법적 분해에는 가약성 개념을 보충부분의 성질로 바꾸어 쓰는 것이 대응한다(\S~3). 그러나 본질적으로 그 공동논문과 같은 고찰을 거치면, 여기에 제시한 모든 정리를 잉여류계와 어떤 부분계의 가법적 분해정리로 볼 수도 있다. 이 잉여류계는 처음에 바탕으로 삼은 환과 똑같이 일반적인 환을 이룬다. 실제로 임의의 환은 환원소들 사이의 모든 항등관계 전체에 대응하는 아이디얼의 잉여류계로 볼 수 있다. 또는 나머지 관계들을 영역 안에서도 성립한다고 가정하면, 이 관계들의 한 부분계에 대응하는 아이디얼의 잉여류계로도 볼 수 있다.

이 논평은 프렝켈의 연구들이 어디에 놓이는지도 보여 준다.\footnote{A. Fraenkel, \emph{Über die Teiler der Null und die Zerlegung von Ringen}. J. f. M. 145 (1914), S. 139. \emph{Über gewisse Teilbereiche und Erweiterungen von Ringen}. Habilitationsschrift, Leipzig, Teubner, 1916. \emph{Über einfache Erweiterungen zerlegbarer Ringe}. J. f. M. 151 (1920), S. 121.} 프렝켈은 네 분해가 일치하도록 제한조건(정칙원소의 존재, 그 원소들에 의한 나눗셈, 분해가능성 조건)을 부과한 환의 가법적 분해를 고찰한다. 네 분해가 일치하기 때문에, 아이디얼이 유한개의 진약수만을 가져야 한다는 그의 유한조건도 --- 그는 일부 경우 이 조건을 쓰지 않기도 한다 --- 우리의 조건보다 더 강한 제한이 아니다. 프렝켈의 출발점은 본질적으로 대수적인 다른 목표들 때문에 정해진 것이며, 그는 대수적 확대를 통해 제한조건이 더 적은 일반적인 환들에 이른다.

\begin{center}
\S~1.\\[0.25em]
\textbf{환영역, 아이디얼, 유한조건.}
\end{center}

\textbf{1.} 바탕 영역 $\Sigma$는 추상적으로 정의된 (가환) \emph{환}이라 하자.\footnote{이 정의는 프렝켈의 교수자격논문에서 제한조건 6, I, II를 빼고 가져온 것이다. 그 대신 덧셈의 교환법칙을 포함해야 했다. 따라서 이는 곱셈의 가역성을 제외한 체의 정의법칙들이다.} 곧 $\Sigma$는 원소 $a,b,c,\ldots,f,g,h,\ldots$의 계이며, 그 안에는 통상적인 조건들을 만족하는 한 관계가 \emph{등식}으로 정의되어 있다. 또 \emph{덧셈}과 \emph{곱셈}이라는 \emph{두} 연산(결합방식)에 의해 임의의 두 환원소 $a,b$에서 각각 합 $a+b$와 곱 $a\cdot b$라는 제3의 원소가 언제나 유일하게 얻어진다. \emph{환}과, 그 밖에는 \emph{완전히} 임의적인 이 \emph{연산들}은 다음 \emph{법칙들}을 만족해야 한다.
\begin{enumerate}
\item \emph{덧셈의 결합법칙: $(a+b)+c=a+(b+c)$.}
\item \emph{덧셈의 교환법칙: $a+b=b+a$.}
\item \emph{곱셈의 결합법칙: $(a\cdot b)\cdot c=a\cdot(b\cdot c)$.}
\item \emph{곱셈의 교환법칙: $a\cdot b=b\cdot a$.}
\item \emph{분배법칙: $a\cdot(b+c)=a\cdot b+a\cdot c$.}
\item \emph{제한 없이 유일한 뺄셈의 법칙.} 방정식 $a+x=b$를 만족하는 원소 $x$가 $\Sigma$에 오직 하나 존재한다. (이를 $x=b-a$로 나타낸다.)
\end{enumerate}
이 성질들에서 영원의 존재가 따른다. 그러나 환에 단위원이 있을 필요는 없으며, 어느 인자도 영이 아닌데 두 원소의 곱이 영이 될 수도 있다. 곱이 영이면 언제나 한 인자가 영이고, 그뿐 아니라 단위원도 갖는 환을 \emph{진정한 정역}이라 한다. 유한합 $a+a+\cdots+a$에는 통상적인 약기호 $na$를 도입한다. 이때 \emph{정수 $n$은 환원소가 아니라 단지 약기호로 보아야 하며}, $a=1\cdot a$, $na+a=(n+1)a$로 재귀적으로 정의된다.

\textbf{2.} $\Sigma$의 한 \emph{아이디얼} $\ideal M$\footnote{아이디얼은 대문자 독일체로 나타낸다. $\ideal M$은 보통 ``모듈'' 또는 형식모듈이라 불리는 다항식 아이디얼의 예를 떠올리게 하려고 택한 기호이다. 덧붙여 \S\S~1--3에서는 아이디얼 성질이 아니라 모듈 성질만을 사용한다. 이에 관해서는 \S~9를 보라.}이란 다음 두 조건을 만족하는 $\Sigma$ 원소들의 계를 뜻한다.
\begin{enumerate}
\item \emph{$\ideal M$이 $f$를 포함하고 $a$가 $\Sigma$의 임의의 원소이면, $\ideal M$은 $a\cdot f$도 포함한다.}
\item \emph{$\ideal M$이 $f,g$를 포함하면 차 $f-g$도 포함한다.} 따라서 $f$와 함께 모든 정수 $n$에 대한 $nf$도 포함한다.
\end{enumerate}
$f$가 $\ideal M$의 원소이면 통상처럼 $f\equiv0(\ideal M)$으로 나타내고, \emph{$f$가 $\ideal M$로 나누어진다}고 말한다. $\ideal N$의 모든 원소가 동시에 $\ideal M$의 원소, 곧 $\ideal M$로 나누어질 때 \emph{$\ideal N$이 $\ideal M$로 나누어진다}고 말하며 $\ideal N\equiv0(\ideal M)$으로 쓴다. $\ideal M$이 $\ideal N$에 없는 원소를 포함하여 역으로 $\ideal N$으로 나누어지지 않으면, $\ideal M$을 $\ideal N$의 \emph{진약수}라 한다. $\ideal N\equiv0(\ideal M)$과 $\ideal M\equiv0(\ideal N)$에서 $\ideal N=\ideal M$이 따른다.

그 밖의 익숙한 개념들도 문자 그대로 유지된다. 두 아이디얼 $\ideal A,\ideal B$의 \emph{최대공약아이디얼} $\ideal D=(\ideal A,\ideal B)$는 $a$가 $\ideal A$의 모든 원소를, $b$가 $\ideal B$의 모든 원소를 돌 때 $a+b$ 꼴로 나타나는 원소 전체를 뜻한다. $\ideal D$는 다시 아이디얼이다. 무한히 많은 아이디얼의 최대공약아이디얼 $\ideal D=(\ideal A_1,\ideal A_2,\ldots,\ideal A_\nu,\ldots)$도 매번 유한개의 아이디얼에서 택한 원소들의 합 $d=a_{i_1}+a_{i_2}+\cdots+a_{i_n}$으로 나타나는 모든 원소 $d$의 전체로 정의된다. 여기서도 $\ideal D$는 다시 아이디얼이다.

특히 아이디얼 $\ideal M$이 유한개의 원소 $f_1,f_2,\ldots,f_\varrho$를 포함하고
\[
 \ideal M=(f_1\ldots f_\varrho);\qquad \text{곧}\qquad
 f=a_1f_1+\cdots+a_\varrho f_\varrho+n_1f_1+\cdots+n_\varrho f_\varrho
\]
가 모든 $f\equiv0(\ideal M)$에 대하여 성립한다고 하자. 여기서 $a_i$들은 \emph{환영역}의 원소이고 $n_i$들은 정수이다. 그러면 $\ideal M$을 \emph{유한아이디얼}, $f_1,\ldots,f_\varrho$를 한 \emph{아이디얼기저}라 한다.

\emph{이하에서는 다음 유한조건을 만족하는 환 $\Sigma$만을 바탕으로 삼는다. $\Sigma$의 각 아이디얼은 유한아이디얼, 곧 아이디얼기저를 갖는 아이디얼이다.}

\textbf{3.} 유한조건에서 이후의 모든 고찰의 바탕이 되는 다음 정리가 바로 따른다.

\textbf{정리 I (유한사슬 정리).}\footnote{데데킨트가 수모듈에 대하여 처음 진술하였다. \emph{Zahlentheorie}, Suppl. XI, \S~172, Satz VIII (4판). 우리의 증명과 ``사슬''이라는 명칭은 그곳에서 가져왔다. 다항식 아이디얼에 대해서는 라스커, 앞의 글 56쪽(보조정리)에 있다. 그러나 두 경우 모두 이 정리는 산발적으로만 쓰였다. 우리의 적용은 모두 선택공리에 바탕을 둔다.} \emph{$\ideal M,\ideal M_1,\ideal M_2,\ldots,\ideal M_\nu,\ldots$가 $\Sigma$의 가산무한 아이디얼계이고 각 아이디얼이 바로 다음 아이디얼로 나누어진다면, 어떤 유한한 지수 $n$부터 모든 아이디얼은 동일하다: $\ideal M_n=\ideal M_{n+1}=\cdots$. 달리 말하면 $\ideal M,\ideal M_1,\ideal M_2,\ldots,\ideal M_s,\ldots$가 각 아이디얼이 바로 앞 아이디얼의 진약수인 단순순서 아이디얼사슬을 이루면, 이 사슬은 유한 단계에서 끝난다.}

실제로 $\ideal D=(\ideal M_1,\ideal M_2,\ldots,\ideal M_\nu,\ldots)$를 이 계의 최대공약아이디얼이라 하고, 유한조건에 의해 언제나 존재하는 $\ideal D$의 기저를 $f_1,\ldots,f_k$라 하자. 나눗셈 가정에 따라 $\ideal D$의 각 원소는 동시에 사슬의 한 아이디얼에 속한다. 실제로
\[
 f=g+h,\qquad g\equiv0(\ideal M_r),\qquad h\equiv0(\ideal M_s),\qquad (r\le s)
\]
이면 $g\equiv0(\ideal M_s)$이고, 따라서 $f\equiv0(\ideal M_s)$이다. $f$가 더 많은 성분의 합일 때도 마찬가지이다. 그러므로 어떤 유한한 지수 $n$이 존재하여
\[
 f_1\equiv0(\ideal M_n);\quad\ldots;\quad f_k\equiv0(\ideal M_n);\quad
 \ideal D=(f_1\ldots f_k)\equiv0(\ideal M_n)
\]
이 된다. 역으로 $\ideal M_n\equiv0(\ideal D)$이므로 $\ideal M_n=\ideal D$이다. 더 나아가
\[
 \ideal M_{n+i}\equiv0(\ideal D);\qquad
 \ideal D=\ideal M_n\equiv0(\ideal M_{n+i})
\]
이므로 모든 $i$에 대하여 $\ideal M_{n+i}=\ideal D=\ideal M_n$이다. 이로써 정리가 증명되었다.

역으로 이 정리에서 다시 아이디얼기저의 존재가 따라온다는 점도 주목하자. 따라서 유한조건을 이렇게 기저를 쓰지 않는 꼴로 진술할 수도 있었다.

\begin{center}
\S~2.\\[0.25em]
\textbf{한 아이디얼을 유한개의 기약아이디얼의 최소공배수로 나타내기.}
\end{center}

아이디얼 $\ideal B_1,\ideal B_2,\ldots,\ideal B_k$의 최소공배수 $[\ideal B_1,\ideal B_2,\ldots,\ideal B_k]$를 통상처럼 $\ideal B_1$로도, $\ideal B_2$로도, \ldots, $\ideal B_k$로도 나누어지는 모든 원소의 전체로 정의하자. 기호로 쓰면
\[
 \text{$f\equiv0(\ideal B_i)$,}\quad (i=1,2,\ldots,k)
 \quad\text{이면}\quad
 f\equiv0([\ideal B_1,\ideal B_2,\ldots,\ideal B_k])
\]
이고, 그 역도 성립한다. 최소공배수는 다시 아이디얼이다. 아이디얼 $\ideal B_i$들을 이 분해의 성분이라고도 한다.

\textbf{정의 I.} \emph{표현 $\ideal M=[\ideal B_1\cdots\ideal B_k]$에서 어느 $\ideal B_i$도 나머지 아이디얼들의 최소공배수 $\ideal A_i$의 약수가 아니고, 어느 $\ideal B_i$도 진약수로 바꿀 수 없으면 이 표현을 축약표현이라 한다.\footnote{축약되지 않은 표현의 한 예는 모든 지수 $\lambda\ge2$에 대한 $(x^2,xy)=[(x),(x^2,xy,y^\lambda)]$이다. $\lambda=1$에 대응하는 표현 $[(x),(x^2,y)]$은 그에 대응하는 축약표현이다. (전쟁에서 전사한 K. 헨첼트가 으뜸아이디얼로의 분해가 유일하지 않음을 보여 주는 가장 간단한 예로 이 표현을 알려 주었다.)} 두 조건이 아이디얼 $\ideal B_i$에 대해서만 만족되면 이 표현을 $\ideal B_i$에 관하여 축약되었다고 한다. 최소공배수 $\ideal A_i=[\ideal B_1,\ldots,\ideal B_{i-1},\ideal B_{i+1},\ldots,\ideal B_k]$를 $\ideal B_i$의 보충부분이라 한다. 첫째 조건만을 만족하는 표현을 최단표현이라 한다.}

이제 한 아이디얼을 최소공배수로 나타낼 때에는 다음 정리에 따라 축약표현만 고찰하면 충분하다.

\textbf{보조정리 I.} \emph{한 아이디얼을 유한개의 아이디얼의 최소공배수로 나타낸 모든 표현은 적어도 한 가지 방식으로 축약표현으로 바꿀 수 있다. 특히 연속적인 분해로 그러한 표현을 얻을 수 있다.}

실제로 $\ideal M=[\ideal B_1^*\cdots\ideal B_l^*]$가 $\ideal M$의 임의의 표현이라고 하자. 남겨 둔 아이디얼들의 최소공배수로 나누어지는 $\ideal B_i^*$들을 차례로 버린다. 나머지 아이디얼들도 여전히 $\ideal M$을 주므로, 이렇게 생긴 표현
\[
 \ideal M=[\ideal B_1\cdots\ideal B_k]=[\ideal A_i,\ideal B_i]
\]
에서는 첫째 조건이 만족되고, 따라서 이 표현은 최단표현이다. 어느 $\ideal B_i$를 진약수로 바꾸어도 이 조건은 유지된다. 둘째 조건은 유한사슬 정리(정리 I)에 따라 언제나 만족시킬 수 있다. 실제로
\[
 \ideal M=[\ideal A_i,\ideal B_i]=[\ideal A_i,\ideal B_i^{(1)}]
 =\cdots=[\ideal A_i,\ideal B_i^{(\nu)}]=\cdots
\]
라고 하자. 여기서 각 $\ideal B_i^{(\nu)}$는 바로 앞 아이디얼의 진약수이다. 그러면 이 정리에 따라 사슬 $\ideal B_i,\ideal B_i^{(1)},\ldots,\ideal B_i^{(\nu)},\ldots$는 유한 단계에서 끝나야 한다. 따라서 표현 $\ideal M=[\ideal A_i,\ideal B_i^{(n)}]$에서는 $\ideal B_i^{(n)}$을 그 진약수로 바꿀 수 없다. $\ideal A_i$를 진약수로 바꾸면 이 사실은 더욱 분명히 유지된다. 그러므로 이미 축약된 $\ideal B$들을 포함하여 매번 보충부분을 만들면서 이 절차를 각 $\ideal B_i$에 차례로 적용하면 축약표현을 얻는다.\footnote{앞의 예는 주어진 표현이 그러한 축약표현을 유일하게 정하지 않음을 보여 준다. $\lambda\ge2$일 때 $(x^2,xy)=[(x),(x^2,xy,y^\lambda)]$에는 $[(x),(x^2,y)]$뿐 아니라 임의의 $\mu$에 대한 $[(x),(x^2,\mu x+y)]$도 대응하는 축약표현이다.}

그러한 표현을 연속적으로 얻으려면, 각각 축약된 표현들
\[
 \ideal M=[\ideal B_1,\ideal C_1];\quad
 \ideal C_1=[\ideal B_2,\ideal C_2];\quad\ldots;\quad
 \ideal C_{l-1}=[\ideal B_l,\ideal C_l]
\]
에서 생기는 표현 $\ideal M=[\ideal B_1\cdots\ideal B_l\ideal C_l]$도 축약되었음을 보이면 된다. 이를 위해서는 축약표현
\[
 \ideal M=[\ideal B_1\cdots\ideal B_l\ideal C];\qquad
 \ideal C=[\ideal C_1,\ideal C_2]
\]
에서 $\ideal M=[\ideal B_1\cdots\ideal B_l\ideal C_1\ideal C_2]$도 축약되었음을 보이면 충분하다. 실제로 가정에 따라 어느 $\ideal B$도 그 보충부분의 약수가 아니다. 어느 $\ideal C_i$가 그 보충부분의 약수라면, 둘째 표현이 축약되어 $\ideal C_1,\ideal C_2$가 $\ideal C$의 진약수라는 사실 때문에 첫째 표현에서 $\ideal C$를 진약수로 바꿀 수 있게 되어 가정에 어긋난다. 따라서 이 표현은 최단이다. 더 나아가 가정에 따라 어느 $\ideal B$도 진약수로 바꿀 수 없다. 어느 $\ideal C_i$를 진약수로 바꿀 수 있다면, $\ideal C$의 표현이 축약되어 있으므로 이는 $\ideal C$를 진약수로 바꾸는 것에 대응하여 다시 가정에 어긋난다. 이로써 보조정리가 증명되었다.

\textbf{정의 II.} \emph{한 아이디얼 $\ideal M$을 두 진약수의 최소공배수로 나타낼 수 있으면 $\ideal M$을 가약이라 하고, 그렇지 않으면 $\ideal M$을 기약이라 한다.}

이제 축약표현을 이용하여 유한사슬 정리인 정리 I로 다음을 증명한다.

\textbf{정리 II.} \emph{각 아이디얼은 유한개의 기약아이디얼의 최소공배수로 나타낼 수 있다.}\footnote{그러한 표현이 일반적으로 유일하지 않다는 것은 임의의 $\mu$에 대한 앞의 예 $(x^2,xy)=[(x),(x^2,\mu x+y)]$가 보여 준다. 두 성분은 임의의 $\mu$에 대하여 기약이다. 실제로 $(x)$의 모든 약수는 $(x,g(y))$ 꼴이며, 여기서 $g(y)$는 $y$의 다항식이다. 따라서 그러한 두 약수의 최소공배수도 이 꼴이고 $(x)$의 진약수가 된다. $(x^2,\mu x+y)$는 유일한 약수 $(x,y)$만을 가지므로 반드시 역시 기약이다.}

실제로 임의의 아이디얼 $\ideal M$은 기약이거나 가약이다. 기약이면 $\ideal M=[\ideal M]$가 정리 II가 요구하는 표현이다. 가약이면 $\ideal M=[\ideal B_1,\ideal C_1]$이고, 여기서 $\ideal B_1,\ideal C_1$은 $\ideal M$의 진약수이며 보조정리 I에 따라 이 표현을 축약되었다고 가정할 수 있다. $\ideal C_1$에도 같은 양자택일이 성립한다. 곧 $\ideal C_1$이 기약이거나, 축약표현 $\ideal C_1=[\ideal B_2,\ideal C_2]$가 존재한다.

계속하면 축약표현들의 열
\srcnumdisplay[3.1em]{(1)}{\begin{aligned}
 \ideal M&=[\ideal B_1,\ideal C_1];\qquad
 \ideal C_1=[\ideal B_2,\ideal C_2];\quad\ldots;\\
 \ideal C_{\nu-1}&=[\ideal B_\nu,\ideal C_\nu];\quad\ldots
\end{aligned}}
을 얻는다. 사슬 $\ideal C_1,\ideal C_2,\ldots,\ideal C_n,\ldots$에서 각 $\ideal C_n$은 바로 앞 아이디얼의 진약수이므로 사슬은 유한 단계에서 끝난다. 따라서 $\ideal C_n$이 기약인 지수 $n$이 존재한다. 보조정리 I에 따라 표현 $\ideal M=[\ideal B_1\cdots\ideal B_n\ideal C_n]=[\ideal A_n,\ideal C_n]$도 축약되어 있다. 따라서 $\ideal C_n$은 보충부분 $\ideal A_n$의 약수가 될 수 없고, 표현 $\ideal M=[\ideal A_n,\ideal C_n]$에서 $\ideal A_n$을 진약수로 바꿀 수 없다. 필요하면 $\ideal A_n$을 진약수로 바꾸어\footnote{실제로 이 표현은 $\ideal A_n$에 관해서도 축약되어 있다. 이는 \S~3의 보조정리 IV에서 보조정리 I의 역으로 증명한다.} 전체 표현을 축약하면, 모든 가약아이디얼이 한 기약아이디얼과 그 보충아이디얼의 최소공배수로 된 축약표현을 가짐을 알 수 있다. 그러므로 (1)에서는 일반성을 잃지 않고 모든 $\ideal B_i$를 기약이라고 가정할 수 있다. 위의 논증을 되풀이하면 기약인 $\ideal C_n$의 존재를 얻고, 이로써 정리 II가 증명된다.

\begin{center}
\S~3.\\[0.25em]
\textbf{서로 다른 두 기약아이디얼 분해에서 성분 수의 동일성.}
\end{center}

성분 수의 동일성을 증명하려면 먼저 한 아이디얼의 가약성 또는 기약성을 그 보충부분의 성질로 나타내야 한다. 이는 다음 정리가 해 준다.

\textbf{정리 III.}\footnote{정리 III은 슈마이들러와 뇌터--슈마이들러의 논문들에서 모듈로부터 잉여군으로 넘어가는 과정에 대응한다(서론 참조). $\ideal C$는 잉여군에, $\ideal N_1$과 $\ideal N_2$는 그 잉여군이 분해되는 부분군들에 대응한다.} \emph{최단표현 $\ideal M=[\ideal A,\ideal C]$가 $\ideal C$에 관하여 축약되어 있다고 하자. 이때 $\ideal C$가 가약이기 위한 필요충분조건은 $\ideal M$의 진약수인 두 아이디얼 $\ideal N_1,\ideal N_2$가 존재하여 다음을 만족하는 것이다.}
\srcnumdisplay[3.1em]{(2)}{\ideal N_1\equiv0(\ideal A);\qquad
 \ideal N_2\equiv0(\ideal A);\qquad
 [\ideal N_1,\ideal N_2]=\ideal M.}
\emph{또한 이로부터 다음이 따른다. 조건 (2)가 성립하고 $\ideal C$가 기약이면, $\ideal N_i$ 중 적어도 하나는 $\ideal M$의 진약수가 아니며 $\ideal N_i=\ideal M$이다.}

$\ideal C=[\ideal C_1,\ideal C_2]$이고 $\ideal C_1,\ideal C_2$가 $\ideal C$의 진약수라고 하자. 그러면
\[
 \ideal M=[\ideal A,\ideal C]=[\ideal A,\ideal C_1,\ideal C_2]
       =[[\ideal A,\ideal C_1],[\ideal A,\ideal C_2]].
\]
여기서 아이디얼 $[\ideal A,\ideal C_i]$는 $\ideal M$의 진약수이다. 그렇지 않으면 $[\ideal A,\ideal C]$가 $\ideal C$에 관하여 축약되어 있지 않기 때문이다. 또한 이 아이디얼들은 $\ideal A$로 나누어지므로 조건 (2)의 필요성이 증명되었다. (표현 (2)는 축약되어 있지 않다. $[\ideal A,\ideal C_i]$를 $\ideal C_i$로 바꿀 수 있기 때문이다.)

이제 거꾸로 (2)가 성립한다고 하자. 아이디얼들을
\[
 \ideal C_1=(\ideal C,\ideal N_1),\qquad
 \ideal C_2=(\ideal C,\ideal N_2),\qquad
 \ideal C^*=[\ideal C_1,\ideal C_2]
\]
로 만든다. 그러면 $\ideal C$는 $\ideal C_1$과 $\ideal C_2$ 각각으로 나누어지고, 따라서 그 최소공배수 $\ideal C^*$로도 나누어진다. $\ideal C^*$가 $\ideal C$로 나누어짐을 보이기 위하여
\[
 f\equiv0(\ideal C^*);\quad f\equiv0(\ideal C_1);\quad f\equiv0(\ideal C_2),
 \quad\text{또는}\quad f=c+n_1=\bar c+n_2
\]
라고 하자. 여기서 $c,\bar c,n_1,n_2$는 각각 $\ideal C,\ideal N_1,\ideal N_2$에 속하는 원소들이며, 특히 $n_i$는 $\ideal N_i$로 나누어진다. 따라서 차
\[
 g=c-\bar c=n_2-n_1
\]
는 $\ideal C$와 $\ideal A$ 양쪽으로 나누어지고, 따라서 $\ideal M$으로 나누어진다. 더 나아가 $n_1=n_2+m$이므로 $n_1$은(마찬가지로 $n_2$도) $\ideal N_1$과 $\ideal N_2$ 양쪽으로 나누어지고, 따라서 $\ideal M$으로 나누어진다. 그러므로
\[
 f=c+m,\qquad f\equiv0(\ideal C),\qquad \ideal C^*=\ideal C
\]
를 얻는다. 이때 $\ideal C_1$과 $\ideal C_2$는 $\ideal C$의 진약수이다. 실제로 $\ideal C_1=(\ideal C,\ideal N_1)=\ideal C$라면 $\ideal N_1$이 $\ideal C$로 나누어질 것이고, $\ideal A$로도 나누어지므로 가정에 어긋나게 $\ideal N_1=\ideal M$이 될 것이다. 따라서 $\ideal C=[\ideal C_1,\ideal C_2]$로부터 $\ideal C$가 가약임을 알 수 있으며, 정리 III이 증명되었다.

거의 같은 논증이 다음 결과도 준다는 점을 덧붙여 둔다.

\textbf{보조정리 II.} \emph{최단표현 $\ideal M=[\ideal A,\ideal C]$에서 아이디얼 $\ideal C$를 진약수로 바꿀 수 있으면 $\ideal C$는 가약이다.}

실제로
\[
 \ideal M=[\ideal A,\ideal C]=[\ideal A,\ideal C_1]
\]
이라고 하고
\[
 \ideal C^*=[\ideal C_1,(\ideal A,\ideal C)]
\]
로 놓자. 그러면 다시 $\ideal C$가 $\ideal C^*$로 나누어진다. 또한 $f\equiv0(\ideal C^*)$로부터
\[
 f=c_1=a+c
\]
를 얻는다. 따라서 차 $a=c_1-c$는 $\ideal A$와 $\ideal C_1$ 양쪽으로 나누어지고, 따라서 $\ideal M$으로 나누어진다. 그러므로 $c_1=c+m$, $f\equiv0(\ideal C)$, $\ideal C^*=\ideal C$이다. 가정에 따라 $\ideal C_1$과 $(\ideal A,\ideal C)$가 모두 $\ideal C$의 진약수이므로, $\ideal C=\ideal C^*$로부터 $\ideal C$가 가약임을 알 수 있다.

따라서 기약아이디얼 $\ideal C$는 진약수로 바꿀 수 없다.

이제 $\ideal M$을 유한개의 기약아이디얼의 최소공배수로 나타내는 서로 다른 두 최단표현
\[
 \ideal M=[\ideal B_1\cdots\ideal B_k]=[\ideal D_1\cdots\ideal D_l]
\]
이 주어졌다고 하자. 보조정리 II 뒤의 주석에 따라 이 표현들은 동시에 축약되어 있다. 먼저 다음을 증명한다.

\textbf{보조정리 III.} \emph{각 보충부분 $\ideal A_i=[\ideal B_1\cdots\ideal B_{i-1}\ideal B_{i+1}\cdots\ideal B_k]$에 대하여 $\ideal M=[\ideal A_i,\ideal D_j]$가 되게 하는 아이디얼 $\ideal D_j$가 존재한다.}

실제로 $\ideal M=[\ideal D_1,\ideal C_1]$, $\ideal C_1=[\ideal D_2,\ideal C_{12}]$ 등으로 놓으면
\[
 \ideal M=[\ideal A_i,\ideal M]=[\ideal A_i,\ideal D_1,\ideal C_1]
 =[[\ideal A_i,\ideal D_1],[\ideal A_i,\ideal C_1]]
\]
을 얻는다. 여기서 $\ideal N_1=[\ideal A_i,\ideal D_1]$, $\ideal N_2=[\ideal A_i,\ideal C_1]$에 대하여 정리 III의 조건 (2)가 성립한다. 실제로 $\ideal M=[\ideal A_i,\ideal B_i]$는 $\ideal B_i$에 관하여 축약되어 있고 최단표현이다. 그런데 $\ideal B_i$가 기약이라고 가정했으므로 반드시 $\ideal N_i$ 중 하나가 $\ideal M$과 같아야 한다.

$\ideal N_1=\ideal M$이면 보조정리가 증명된다. $\ideal N_2=\ideal M$이면 그에 따라 $\ideal M=[[\ideal A_i,\ideal D_2],[\ideal A_i,\ideal C_{12}]]$가 되고, 같은 논증에 의하여 여기서도 한 성분이 $\ideal M$과 같아야 한다. 이 과정을 계속하면 $j<l$인 어떤 $j$에 대하여 $\ideal M=[\ideal A_i,\ideal D_j]$가 되거나, 아니면 $\ideal M=[\ideal A_i,\ideal C_{1\ldots l-1}]$가 된다. $\ideal C_{1\ldots l-1}=\ideal D_l$이므로 후자의 경우에도 보조정리가 증명된다.

이제 이로부터 다음을 얻는다.

\textbf{정리 IV.} \emph{한 아이디얼을 기약아이디얼들의 최소공배수로 나타내는 서로 다른 두 최단표현에서 성분의 수는 같다.}

실제로 $i=1$에 대하여 보조정리로부터
\[
 \ideal M=[\ideal A_1,\ideal B_1]=[\ideal A_1,\ideal D_{j_1}]
 = [\ideal D_{j_1},\ideal B_2,\ldots,\ideal B_k]
\]
를 얻는다. 이제 두 분해
\[
 \ideal M=[\ideal D_{j_1},\ideal B_2,\ldots,\ideal B_k]
          =[\ideal D_1,\ideal D_2,\ldots,\ideal D_l]
\]
를 생각하고, $\ideal B_2$의 보충부분 $\bar{\ideal A}_2=[\ideal D_{j_1},\ideal B_3,\ldots,\ideal B_k]$에 관하여 위의 논증을 반복하면
\[
 \ideal M=[\bar{\ideal A}_2,\ideal B_2]=[\bar{\ideal A}_2,\ideal D_{j_2}]
           =[\ideal D_{j_1},\ideal D_{j_2},\ideal B_3,\ldots,\ideal B_k]
\]
를 얻으며, 이 절차를 계속하면
\[
 \ideal M=[\ideal D_{j_1},\ideal D_{j_2},\ldots,\ideal D_{j_k}]
\]
가 된다. 가정에 따라 표현 $\ideal M=[\ideal D_1\cdots\ideal D_l]$은 최단이므로 어느 $\ideal D$도 생략할 수 없다. 따라서 $\ideal D_{j_i}$들 가운데 서로 다른 것들이 모든 $\ideal D$를 빠짐없이 포함해야 하며, 이에 따라 $k\ge l$이다. 보조정리와 이어지는 모든 논증에서 $\ideal B$와 $\ideal D$를 서로 바꾸면 마찬가지로 $l\ge k$를 얻는다. 그러므로 $k=l$이고, 이로써 성분 수의 동일성이 증명되었다. 또한 이로부터 아이디얼 $\ideal D_{j_i}$들이 모두 서로 다름을 알 수 있다. 그렇지 않으면 $\ideal D_{j_i}$들로 된 최단표현에 $k$개보다 적은 성분이 나타날 것이기 때문이다. 따라서 지표를 $\ideal D_{j_i}=\ideal D_i$가 되도록 정할 수 있다. 같은 이유로 모든 중간표현 $\ideal M=[\ideal D_{j_1}\cdots\ideal D_{j_r},\ideal B_{r+1},\ldots,\ideal B_k]$도 최단이고, 따라서 보조정리 II 뒤의 주석에 따라 축약되어 있다.

성분 수의 동일성은 다음 보조정리를 통하여 보조정리 I의 역을 준다.

\textbf{보조정리 IV.} \emph{축약표현의 성분들을 여러 묶음으로 모으고 각 묶음의 최소공배수를 취하면, 그 결과로 생기는 표현도 축약되어 있다. 다시 말해 축약표현}
\[
 \ideal M=[\ideal C_{11}\cdots\ideal C_{1\mu_1};\;\ldots;\;
 \ideal C_{\sigma1}\cdots\ideal C_{\sigma\mu_\sigma}]
\]
\emph{에서 $\ideal N_i=[\ideal C_{i1}\cdots\ideal C_{i\mu_i}]$로 놓으면 $\ideal M=[\ideal N_1\cdots\ideal N_\sigma]=[\ideal N_i,\ideal L_i]$도 축약되어 있다.}

먼저 $\ideal N_i$는 그 보충부분 $\ideal L_i$의 약수가 될 수 없다. 실제로 그 약수 $\ideal C_{ij}$ 가운데 어느 것도 $\ideal L_i$의 약수가 아니기 때문이다. 따라서 이 표현은 최단이다. $\ideal N_i$를 진약수로 바꿀 수 없음을 보이기 위하여 각 $\ideal C$를 그 기약아이디얼 $\ideal B$들로 분해한다.\footnote{여기서는 언제나 $\ideal B$들로 된 최단표현, 따라서 축약표현을 뜻한다.} 그러면 보조정리 I에 따라 축약된 표현들
\[
 \ideal M=[\ideal B_{11}\cdots\ideal B_{1\lambda_1};\;\ldots;\;
 \ideal B_{\sigma1}\cdots\ideal B_{\sigma\lambda_\sigma}];
 \qquad
 \ideal N_i=[\ideal B_{i1}\cdots\ideal B_{i\lambda_i}]
\]
을 얻는다. 이제 $\ideal M=[\ideal N_i^*,\ideal L_i]$가 $\ideal N_i^*$에 관하여 축약되어 있고, $\ideal N_i^*$가 $\ideal N_i$의 진약수라고 하자. 보조정리 II에 따라
\[
 \ideal N_i=[\ideal N_i^*,(\ideal N_i,\ideal L_i)]
\]
이다. 이 표현은 반드시 $\ideal N_i^*$에 관하여 축약되어 있다. 그렇지 않으면 $\ideal M$에서도 $\ideal N_i^*$를 진약수로 바꿀 수 있기 때문이다. 이제 필요하면 $(\ideal N_i,\ideal L_i)$도 진약수로 바꾸어 $\ideal N_i$의 축약표현을 얻는다. $\ideal N_i$의 이 두 성분을 기약아이디얼들로 분해하면, $\ideal N_i$의 기약아이디얼 수 $\lambda_i$는 두 성분의 기약아이디얼 수의 합이 된다. 따라서 진약수 $\ideal N_i^*$에 대응하는 기약아이디얼의 수는 반드시 $\lambda_i$보다 작다. 그러면 $\ideal M=[\ideal N_i^*,\ideal L_i]$를 기약아이디얼들로 분해할 때에도 아이디얼의 수가 $\sum_i\lambda_i$보다 작아지므로, 성분 수의 동일성에 모순된다. 특수한 경우 $\sigma=2$로부터 표현 $\ideal M=[\ideal N_i,\ideal L_i]$가 보충부분 $\ideal L_i$에 관해서도 축약되어 있음이 또한 따른다.

\begin{center}
\S~4.\\[0.25em]
\textbf{으뜸아이디얼. 서로 다른 두 기약아이디얼 분해에서 대응 소아이디얼의 유일성.}
\end{center}

이하에서는 으뜸아이디얼과 기약아이디얼 사이의 관계를 다룬다.

\textbf{정의 III.} \emph{$a\cdot b\equiv0(\ideal Q)$, $a\not\equiv0(\ideal Q)$로부터 반드시 $b^\varkappa\equiv0(\ideal Q)$가 따르고 지수 $\varkappa$가 유한한 수이면, 아이디얼 $\ideal Q$를 으뜸아이디얼이라 한다.}

이 정의는 다음과 같이 말할 수도 있다. 곱 $a\cdot b$가 $\ideal Q$로 나누어지면 한 인자가 $\ideal Q$로 나누어지거나, 각 인자의 어떤 거듭제곱이 $\ideal Q$로 나누어진다. \emph{특히 $\varkappa$를 언제나 1로 취할 수 있으면 그 아이디얼을 소아이디얼이라 한다.}

으뜸아이디얼(또는 소아이디얼)의 정의와 기저의 존재로부터, 아이디얼의 곱만을 사용하는 다음 정의를 얻는다.\footnote{두 아이디얼의 곱 $\ideal A\cdot\ideal B$는 통상적으로 원소 $a\cdot b$ 전체와 그 유한합들로 이루어진 아이디얼을 뜻한다.}

\textbf{정의 IIIa.} \emph{$\ideal A\cdot\ideal B\equiv0(\ideal Q)$, $\ideal A\not\equiv0(\ideal Q)$로부터 반드시 $\ideal B^\lambda\equiv0(\ideal Q)$가 따르면 아이디얼 $\ideal Q$를 으뜸아이디얼이라 한다. $\lambda$를 언제나 1로 취할 수 있으면 그 아이디얼을 소아이디얼이라 한다. 따라서 소아이디얼 $\ideal P$에 대해서는 $\ideal A\cdot\ideal B\equiv0(\ideal P)$, $\ideal A\not\equiv0(\ideal P)$로부터 언제나 $\ideal B\equiv0(\ideal P)$가 따른다.}

IIIa에서 $\ideal A=(a)$, $\ideal B=(b)$로 놓으면 정의 III이 특수한 경우로 포함되므로, IIIa에 따라 으뜸인 모든 아이디얼은 III에 따라서도 으뜸이다. 거꾸로 $\ideal Q$가 III에 따라 으뜸이고 IIIa의 가정 $\ideal A\cdot\ideal B\equiv0(\ideal Q)$가 성립한다고 하자. 그러면 $\ideal A\equiv0(\ideal Q)$이거나, 아니면 $a\cdot\ideal B\equiv0(\ideal Q)$이면서 $a\not\equiv0(\ideal Q)$인 원소 $a\equiv0(\ideal A)$가 적어도 하나 존재한다. 이제 $b_1,\ldots,b_r$가 $\ideal B$의 아이디얼기저이면 $a\cdot b_i\equiv0(\ideal Q)$이므로 정의 III에 따라
\[
 b_1^{\varkappa_1}\equiv0(\ideal Q);\quad\ldots;\quad
 b_r^{\varkappa_r}\equiv0(\ideal Q)
\]
를 얻는다. 또한
\[
 b=f_1b_1+\cdots+f_r b_r+n_1b_1+\cdots+n_r b_r
\]
이므로 $\lambda=\varkappa_1+\cdots+\varkappa_r$로 놓으면 $\ideal B$의 원소 $\lambda$개로 이루어진 임의의 곱이 $\ideal Q$로 나누어진다. 이로써 III에 따라 으뜸인 아이디얼이 정의 IIIa도 만족함이 증명되었다. 특히 소아이디얼 $\ideal P$의 경우에는 $a\cdot\ideal B\equiv0(\ideal P)$, 따라서 모든 $b\equiv0(\ideal B)$에 대하여 $a\cdot b\equiv0(\ideal P)$이고 $a\not\equiv0(\ideal P)$이면 $b\equiv0(\ideal P)$, 따라서 $\ideal B\equiv0(\ideal P)$가 따른다. 이로써 두 정의의 동치가 증명되었다.

으뜸아이디얼과 소아이디얼의 관계는 다음 관찰로 얻어진다. 어떤 거듭제곱이 $\ideal Q$로 나누어지는 성질을 가진 모든 원소 $p$의 전체 $\ideal P$는 소아이디얼을 이룬다. 먼저 $p_1,p_2$와 함께 $a p_1$과 $p_1-p_2$도 이 성질을 가지므로 $\ideal P$가 아이디얼임은 분명하다. 또한 정의 IIIa에서 사용한 기저 논증에 따라 어떤 수 $\lambda$가 존재하여 $\ideal P^\lambda\equiv0(\ideal Q)$가 된다.

이제
\[
 a b\equiv0(\ideal P),\qquad a\not\equiv0(\ideal P)
\]
라고 하자. $\ideal P$의 정의에 따라
\[
 a^\lambda b^\lambda\equiv0(\ideal Q),\qquad a^\lambda\not\equiv0(\ideal Q)
\]
이고, 따라서 $\ideal Q$의 정의에 따라
\[
 b^{\lambda\varkappa}\equiv0(\ideal Q),\qquad\text{따라서}\qquad b\equiv0(\ideal P)
\]
이다. 이로써 $\ideal P$가 소아이디얼임이 증명되었다. 또한 $\ideal P$는 어떤 거듭제곱이 $\ideal Q$로 나누어지는 성질을 가진 모든 아이디얼 $\ideal B$의 최대공약아이디얼로 정의할 수 있다. 실제로 그러한 각 $\ideal B$는 $\ideal P$로 나누어지므로, 이 $\ideal B$들의 최대공약아이디얼 $\ideal D$도 $\ideal P$로 나누어진다. 거꾸로 $\ideal P$ 자체가 그러한 아이디얼 $\ideal B$이므로 $\ideal D$로 나누어지고, 따라서 $\ideal P=\ideal D$이다. 그러므로 $\ideal P$는 $\ideal Q$의 약수이고 그 어떤 거듭제곱이 $\ideal Q$로 나누어지는 소아이디얼이며, 이 성질로 유일하게 정해진다. 실제로
\[
 \ideal Q\equiv0(\ideal P),\qquad \ideal P^\lambda\equiv0(\ideal Q),
 \qquad
 \ideal Q\equiv0(\bar{\ideal P}),\qquad \bar{\ideal P}^{\mu}\equiv0(\ideal Q)
\]
이면
\[
 \ideal P^\lambda\equiv0(\bar{\ideal P}),\qquad
 \bar{\ideal P}^{\mu}\equiv0(\ideal P)
\]
이고, 소아이디얼의 성질에 따라
\[
 \ideal P\equiv0(\bar{\ideal P}),\qquad
 \bar{\ideal P}\equiv0(\ideal P),\qquad
 \ideal P=\bar{\ideal P}
\]
를 얻는다. 정리하면 다음과 같다.

\textbf{정리 V.} \emph{각 으뜸아이디얼 $\ideal Q$에 대하여, $\ideal Q$의 약수이면서 그 어떤 거듭제곱이 $\ideal Q$로 나누어지는 소아이디얼 $\ideal P$가 유일하게 존재한다. $\ideal P$를 $\ideal Q$에 ``대응 소아이디얼''이라 한다.}\footnote{역이 성립하지 않음은 예 $\ideal M=(x^2,xy)$가 보여 준다. 소아이디얼 $(x)$는 모든 조건을 만족하지만 $\ideal M$은 으뜸아이디얼이 아니다.} \emph{$\ideal P$는 어떤 거듭제곱이 $\ideal Q$로 나누어지는 성질을 가진 모든 아이디얼 $\ideal B$의 최대공약아이디얼로 정의된다. $\ideal P^\varrho\equiv0(\ideal Q)$가 되게 하는 가장 작은 수를 $\varrho$라 하면, $\varrho$를 $\ideal Q$의 지수라 한다.}\footnote{유리정수 또는 대수적 정수의 환영역에서와 달리 일반적으로 $\ideal P^\varrho=\ideal Q$가 되지 않음은 다음 예가 보여 준다. $\ideal Q=(x^2,y)$, $\ideal P=(x,y)$, $\ideal P^2=(x^2,xy,y^2)\equiv0(\ideal Q)$이지만 $\ideal Q\not\equiv0(\ideal P^2)$이다. 따라서 $\ideal Q$는 $\ideal P^2$와 다르다.}

이제 으뜸성과 기약성 사이의 관계로서 다음을 증명한다.

\textbf{정리 VI.} \emph{으뜸아이디얼이 아닌 모든 아이디얼은 가약이다. 다시 말해 모든 기약아이디얼은 으뜸아이디얼이다.}\footnote{역이 성립하지 않음은 예를 들어 $\lambda\ge2$일 때 $\ideal Q=(x^2,xy,y^\lambda)=[(x^2,y),(x,y^\lambda)]$가 보여 준다. 여기서 $\ideal Q$는 으뜸이지만 가약이다. $\ideal Q$가 으뜸임은 이 아이디얼이 $x,y$의 모든 $\lambda$차 멱곱을 포함한다는 데서 따른다. 따라서 상수항이 없는 모든 다항식의 어떤 거듭제곱은 $\ideal Q$로 나누어진다. 그런데 $a\cdot b\equiv0(\ideal Q)$에서 다항식 $b$가 상수항을 포함한다면, 즉 모든 $\varkappa$에 대하여 $b^\varkappa\not\equiv0(\ideal Q)$이면, $\ideal Q$의 기저다항식들이 동차이므로 $a\cdot b$의 각 동차성분이 $\ideal Q$로 나누어진다는 사실에서 $a$가 $\ideal Q$로 나누어져야 한다.}

$\ideal K$가 으뜸이 아닌 아이디얼이라고 하자. 그러면 정의 III에 따라 적어도 한 원소쌍 $a,b$가 존재하여
\srcnumdisplay[3.1em]{(3)}{a\cdot b\equiv0(\ideal K);\qquad
 a\not\equiv0(\ideal K);\qquad
 b^\varkappa\not\equiv0(\ideal K)\quad\text{모든 }\varkappa\text{에 대하여}}
를 만족한다. 이제 두 아이디얼
\[
 \ideal L_0=(\ideal K,a);\qquad \ideal N_0=(\ideal K,b)
\]
를 만든다. 이들은 (3)에 따라 $\ideal K$의 진약수이며, 또한 (3)에 따라
\srcnumdisplay[3.1em]{(4)}{\ideal L_0\cdot\ideal N_0\equiv0(\ideal K)}
이다. 최소공배수 $\ideal K_0=[\ideal L_0,\ideal N_0]$의 원소 $f$에 대해서는 다음 둘 중 하나가 성립한다.

\emph{첫째,}
\[
 f\equiv0(\ideal L_0);\qquad f\equiv0(\ideal N_0),
 \qquad\text{즉}\qquad f\equiv a_1\cdot b\;(\ideal K)
\]
로부터 언제나
\[
 f\equiv l_0\cdot b\;(\ideal K);\qquad l_0\equiv0(\ideal L_0)
\]
인 표현이 따른다. 그러면 (4)에 따라 $f\equiv0(\ideal K)$이고, 따라서 $\ideal K_0\equiv0(\ideal K)$이다. 또한 $\ideal K\equiv0(\ideal K_0)$이므로 $\ideal K=\ideal K_0$이다. 이로써 $\ideal K$가 가약임이 증명된다.

\emph{둘째,} 그러한 $l_0$가 존재하지 않는 $f\equiv0(\ideal K_0)$가 적어도 하나 있다. 이때 이 $f$에 대응하는 $a_1$을 이용하여
\[
 \ideal L_1=(\ideal L_0,a_1)=(\ideal K,a,a_1);\qquad
 \ideal N_1=(\ideal K,b^2)
\]
를 만든다. $a_1\cdot b\equiv0(\ideal L_0)$이므로 (4)에 따라
\srcnumdisplay[3.1em]{(4')}{\ideal L_1\cdot\ideal N_1\equiv0(\ideal K)}
도 성립하고, $\ideal L_1$은 $\ideal L_0$의 진약수가 된다.

$\ideal K_1=[\ideal L_1,\ideal N_1]$의 원소 $f$에 대해서도 같은 양자택일이 성립한다.

\emph{첫째,}
\[
 f\equiv0(\ideal L_1);\qquad f\equiv0(\ideal N_1),
\]
즉
\[
 f\equiv a_2\cdot b^2\;(\ideal K)
 \quad\text{이면 언제나}\quad
 f\equiv l_1\cdot b^2;\qquad l_1\equiv0(\ideal L_1)
\]
가 따른다. 그러면 (4')에 따라 $\ideal K=\ideal K_1$이다.

\emph{둘째,} 적어도 한 $f$에 대하여 그러한 $l_1$이 존재하지 않는다. 이 경우 $\ideal L_2=(\ideal L_1,a_2)$, $\ideal N_2=(\ideal K,b^{2^2})$를 만들게 되며, $\ideal L_2\cdot\ideal N_2\equiv0(\ideal K)$이고 $\ideal L_2$는 $\ideal L_1$의 진약수이다. 이 과정을 계속하여 일반적으로
\[
 \ideal L_0=(\ideal K,a);\quad
 \ideal L_1=(\ideal L_0,a_1);\quad\ldots;\quad
 \ideal L_\nu=(\ideal L_{\nu-1},a_\nu);\quad\ldots,
\]
\[
 \ideal N_0=(\ideal K,b);\quad
 \ideal N_1=(\ideal K,b^2);\quad\ldots;\quad
 \ideal N_\nu=(\ideal K,b^{2^\nu});\quad\ldots
\]
로 정의한다. 여기서 $a_i$들은 어떤 $f$가 존재하여
\[
 f\equiv0(\ideal L_{i-1});\qquad f\equiv0(\ideal N_{i-1}),
\]
즉
\[
 f\equiv a_i\cdot b^{2^{i-1}}\;(\ideal K);
 \qquad\text{그러나}\qquad a_i\not\equiv0(\ideal L_{i-1})
\]
가 되도록 정해진다. 그러면 일반적으로 $\ideal L_i\cdot\ideal N_i\equiv0(\ideal K)$이고, (3)에 따라 $\ideal N_i$는 $\ideal K$의 진약수이며 $\ideal L_i$는 $\ideal L_{i-1}$의 진약수이다. 따라서 유한사슬 정리인 정리 I에 따라 $\ideal L_i$들의 사슬은 이를테면 $\ideal L_n$에서 유한하게 끝나야 한다. 그러면 모든 $f\equiv0(\ideal L_n)$, $f\equiv0(\ideal N_n)$에 대하여 $l_n\equiv0(\ideal L_n)$인 $l_n$이 존재하여 $f\equiv l_n\cdot b^{2^n}\;(\ideal K)$가 된다. 따라서 위의 논증에 따라 $\ideal K=[\ideal L_n,\ideal N_n]$이며, 이로써 $\ideal K$가 가약임이 증명되었다.

방금 증명한 결과로부터 \emph{대응 소아이디얼의 유일성}을 다음과 같이 얻는다.

\[
 \ideal M=[\ideal B_1\cdots\ideal B_k]=[\ideal D_1\cdots\ideal D_k]
\]
가 $\ideal M$을 기약아이디얼들의 최소공배수로 나타내는 두 최단표현이라고 하자. 이들은 축약되어 있으며, 정리 IV에 따라 성분의 수도 같다. 또한 정리 IV에 따라 거기서 나타난 중간표현들은(그곳에서 언급했듯 지표를 $j_i=i$로 놓을 수 있다)
\[
 \ideal M=[\ideal D_1\cdots\ideal D_{i-1}\ideal B_i\ideal B_{i+1}\cdots\ideal B_k]
       =[\ideal D_1\cdots\ideal D_{i-1}\ideal D_i\ideal B_{i+1}\cdots\ideal B_k]
       =[\bar{\ideal U}_i,\ideal B_i]=[\bar{\ideal U}_i,\ideal D_i]
\]
처럼 최단표현이다. 따라서
\[
 \bar{\ideal U}_i\cdot\ideal B_i\equiv0(\ideal D_i),\qquad
 \bar{\ideal U}_i\not\equiv0(\ideal D_i),
 \qquad
 \bar{\ideal U}_i\cdot\ideal D_i\equiv0(\ideal B_i),\qquad
 \bar{\ideal U}_i\not\equiv0(\ideal B_i)
\]
이다. 이제 정리 VI에 따라 기약아이디얼 $\ideal B_i$와 $\ideal D_i$는 으뜸이므로, 두 수 $\lambda_i,\mu_i$가 존재하여
\srcnumdisplay[3.1em]{(5)}{\ideal B_i^{\lambda_i}\equiv0(\ideal D_i);\qquad
 \ideal D_i^{\mu_i}\equiv0(\ideal B_i)}
가 된다. $\ideal B_i$와 $\ideal D_i$에 대응 소아이디얼을 각각 $\ideal P_i$, $\bar{\ideal P}_i$라 하자. 곧 $\ideal P_i^{\varrho_i}\equiv0(\ideal B_i)$, $\bar{\ideal P}_i^{\sigma_i}\equiv0(\ideal D_i)$라 하자. 그러면 (5)에 따라
\[
 \ideal P_i^{\lambda_i\varrho_i}\equiv0(\bar{\ideal P}_i),
 \qquad
 \bar{\ideal P}_i^{\mu_i\sigma_i}\equiv0(\ideal P_i)
\]
이고, 소아이디얼의 성질에 따라
\[
 \ideal P_i\equiv0(\bar{\ideal P}_i),
 \qquad
 \bar{\ideal P}_i\equiv0(\ideal P_i),
 \qquad
 \ideal P_i=\bar{\ideal P}_i
\]
를 얻는다. 이로써 다음이 증명되었다.

\textbf{정리 VII.} \emph{한 아이디얼을 기약아이디얼들의 최소공배수로 나타내는 서로 다른 두 최단표현에서 대응 소아이디얼들은 서로 일치한다. 그 가운데 같은 소아이디얼이 나타날 수도 있으며, 각 분해에서 같은 횟수로 나타난다.\footnote{예를 들어 각주 19의 예 $\lambda\ge2$일 때 $(x^2,xy,y^\lambda)=[(x^2,y),(x,y^\lambda)]$가 이를 보여 준다. 두 성분에 대응 소아이디얼은 모두 $(x,y)$이다.} 따라서 두 표현의 아이디얼들은 적어도 한 가지 방식으로 쌍을 지어, 각 아이디얼 $\ideal B_i$의 어떤 거듭제곱이 짝지어진 $\ideal D_i$로 나누어지고 그 역도 성립하게 할 수 있다. 그 수는 정리 IV에 따라 같다.}\footnote{기약아이디얼들 가운데 들어 있는 ``고립'' 아이디얼의 유일성에 대해서는 \S~7을 참조하라.}

\begin{center}
\S~5.\\[0.25em]
\textbf{한 아이디얼을 극대 으뜸아이디얼들의 최소공배수로 나타내기. 대응 소아이디얼의 유일성.}
\end{center}

\textbf{정의 IV.} \emph{최단표현 $\ideal M=[\ideal Q_1\cdots\ideal Q_\alpha]$에서 모든 $\ideal Q$가 으뜸이지만 어느 두 $\ideal Q$의 최소공배수도 더 이상 으뜸이 아니면, 이 표현을 극대 으뜸아이디얼들의 최소공배수라고 한다.}

그러한 표현이 적어도 하나 언제나 존재한다는 사실은 $\ideal M$을 기약아이디얼들의 최소공배수로 나타내는 표현에서 따른다. 실제로 기약아이디얼들은 으뜸이다. 이미 극대 으뜸아이디얼들로 된 표현이거나, 아니면 어떤 두 아이디얼의 최소공배수가 다시 으뜸이다. 후자의 경우 아이디얼의 수가 하나 줄어들므로, 이 절차를 반복하면 유한 단계 뒤에 원하는 표현을 얻는다.

이 표현은 보조정리 IV에 따라 축약되어 있다. 거꾸로 극대 으뜸아이디얼들로 된 모든 축약표현은 $\ideal Q$들을 기약아이디얼들로 분해해 보면 이러한 방식으로 생긴다.

여기서 정리 VII로부터 그에 대응하는 유일성 정리를 얻으려면 대응 소아이디얼과의 관계를 다음과 같이 조사해야 한다.

\textbf{정리 VIII.} \emph{으뜸아이디얼 $\ideal N_1,\ideal N_2,\ldots,\ideal N_\lambda$가 모두 같은 대응 소아이디얼 $\ideal P$를 가지면, 그 최소공배수 $\ideal Q=[\ideal N_1,\ideal N_2,\ldots,\ideal N_\lambda]$도 으뜸이고 $\ideal P$를 대응 소아이디얼로 가진다. 거꾸로 $\ideal Q=[\ideal N_1\cdots\ideal N_\lambda]$가 으뜸아이디얼 $\ideal Q$의 축약표현이면, 모든 $\ideal N_i$는 으뜸이고 $\ideal Q$에 대응 소아이디얼 $\ideal P$를 자신의 대응 소아이디얼로 가진다.}

주장의 첫 부분을 증명하기 위해 먼저 각 $i$에 대하여 $\ideal P^{\varrho_i}\equiv0(\ideal N_i)$이면 $\ideal P^\tau\equiv0(\ideal Q)$도 성립함을 관찰하자. 여기서 $\tau$는 지수 $\varrho_i$들 가운데 가장 큰 것이다. 또한 $\ideal P$는 $\ideal Q$의 약수이므로, $\ideal Q$가 으뜸이기만 하면 $\ideal P$는 반드시 그 대응 소아이디얼이다. 이제
\[
 \ideal A\cdot\ideal B\equiv0(\ideal Q),\qquad
 \ideal B^k\not\equiv0(\ideal Q)\quad\text{(모든 }k\text{에 대하여)}
\]
이면
\[
 \ideal B\not\equiv0(\ideal P);
 \quad\text{따라서}\quad
 \ideal B^k\not\equiv0(\ideal N_i)\quad\text{(모든 }k\text{에 대하여)}
\]
이다. 그러므로 $\ideal A\equiv0(\ideal N_i)$이고, 따라서 $\ideal A\equiv0(\ideal Q)$이다. 이로써 $\ideal Q$가 으뜸이고 $\ideal P$가 그 대응 소아이디얼임이 증명되었다.

이제 거꾸로 먼저
\[
 \ideal Q=[\ideal N_1\cdots\ideal N_\lambda]=[\ideal N_i,\ideal L_i]
\]
가 으뜸아이디얼 $\ideal N_i$들로 된 $\ideal Q$의 최단표현이고, 각 $\ideal P_i$가 $\ideal N_i$에 대응 소아이디얼이라고 하자. 그러면
\[
 \ideal L_i\cdot\ideal N_i\equiv0(\ideal Q),\qquad
 \ideal L_i\not\equiv0(\ideal Q)
\]
(최단표현이므로)에서
\[
 \ideal N_i^{\sigma_i}\equiv0(\ideal Q),\qquad
 \text{따라서}\qquad
 \ideal P_i^{\varrho_i\sigma_i}\equiv0(\ideal Q)
\]
를 얻는다. 동시에 $\ideal P_i$는 $\ideal Q$의 약수이므로, 각 $i$에 대하여 $\ideal P_i$는 $\ideal Q$에 대응 소아이디얼 $\ideal P$와 일치한다.

이제 모든 축약표현 $\ideal Q=[\ideal N_1\cdots\ideal N_\lambda]$에서 $\ideal N_i$들이 으뜸임을 보이는 일만 남았다.\footnote{여기서 축약표현이라는 가정이 본질적임은 축약되지 않은 다음 표현이 보여 준다. $\lambda\ge2$일 때 $\ideal Q=[x^2,xy,y^\lambda]=[(x^2,xy,y^2,yz),(x,y^\lambda)]$이다. 여기서 $(x^2,xy,y^2,yz)=[(x^2,y),(x,y^2,z)]$는 위에서 증명한 바에 따라 으뜸이 아니다. 실제로 뒤의 표현은 으뜸아이디얼들로 된 최단표현이지만 대응 소아이디얼 $(x,y)$와 $(x,y,z)$가 서로 다르다. ($\ideal Q$가 으뜸임은 각주 19에 따른다.)} 이를 위해 $\ideal N_i$들을 그 기약아이디얼 $\ideal B$들로 분해하자. 이렇게 생긴 최단표현 $\ideal Q=[\ideal B_1\cdots\ideal B_s]$에서는 각 아이디얼이 으뜸이고, 방금 증명한 결과에 따라 $\ideal P$를 대응 소아이디얼로 가진다. 그러면 위에서 증명한 정리의 첫 부분에 따라 각 $\ideal N_i$에도 같은 사실이 성립한다. 이로써 정리 VIII이 완전히 증명되었다.

\textbf{따름정리.} 이로부터 소아이디얼은 반드시 기약임이 또한 따른다. 실제로 축약표현 $\ideal P=[\ideal N,\ideal K]$에서 정리 VIII에 따라 $\ideal N\equiv0(\ideal P)$이고, $\ideal P\equiv0(\ideal N)$이기도 하므로 $\ideal P=\ideal N$이며 마찬가지로 $\ideal P=\ideal K$이다. $\ideal P$의 기약성은 직접 얻을 수도 있다. 실제로 $\ideal P=[\ideal N,\ideal K]$이면 $\ideal N\cdot\ideal K\equiv0(\ideal P)$이지만 $\ideal N\not\equiv0(\ideal P)$, $\ideal K\not\equiv0(\ideal P)$여서 소아이디얼의 정의 성질에 모순된다.

이제
\[
 \ideal M=[\ideal Q_1\cdots\ideal Q_\alpha]
          =[\bar{\ideal Q}_1\cdots\bar{\ideal Q}_\beta]
\]
가 $\ideal M$을 극대 으뜸아이디얼들의 최소공배수로 나타내는 두 축약표현이라고 하자. $\ideal Q$들을 기약아이디얼 $\ideal B$들로, $\bar{\ideal Q}$들을 그에 대응하여 기약아이디얼 $\ideal D$들로 분해하면, $\ideal M$을 기약아이디얼들의 최소공배수로 나타내는 두 축약표현을 얻는다. 정리 VII에 따라 두 표현에서는 성분의 수와 대응 소아이디얼들이 모두 일치한다. 정리 VIII에 따라 고정된 $\ideal Q_i$에서 나타나는 모든 기약아이디얼 $\ideal B$는 같은 대응 소아이디얼 $\ideal P_i$를 가진다. 한편 $\ideal Q_j$에 속하는 $\ideal P_j$는 반드시 $\ideal P_i$와 다르다. 그렇지 않으면 정리 VIII에 따라 극대 으뜸아이디얼들로 된 표현이 아니기 때문이다. 따라서 $\ideal Q$들의 수 $\alpha$는 $\ideal B$들에 결부된 서로 다른 소아이디얼 $\ideal P$들의 수와 같고, 이 서로 다른 $\ideal P$들이 $\ideal Q$들에 대응 소아이디얼들이다. $\bar{\ideal Q}$들을 $\ideal D$들로 분해한 것에 대해서도 같은 사실이 성립한다. 그러므로 정리 VII로부터 $\ideal Q$와 $\bar{\ideal Q}$의 수가 같고 그 대응 소아이디얼들이 일치함을 얻는다. 동시에 이 절의 첫머리에서 설명한, 기약아이디얼들을 극대 으뜸아이디얼들로 모으는 과정은 같은 대응 소아이디얼을 가진 것들을 모두, 그리고 오직 그것들만 한데 모으는 과정임을 알 수 있다. 더 나아가 정리 VIII은 극대 으뜸아이디얼의 기약성 성질을 보여 준다. 즉 극대 으뜸아이디얼은 극대 으뜸아이디얼들의 최소공배수로 된 축약표현을 허용하지 않는다.

정리하면 다음이 증명되었다.

\textbf{정리 IX.} \emph{한 아이디얼을 극대 으뜸아이디얼들의 최소공배수로 나타내는 두 축약표현에서 성분의 수와 대응 소아이디얼들이 일치하며, 이 소아이디얼들은 모두 서로 다르다. 다시 말해 각 $\ideal Q$에 $\bar{\ideal Q}$ 하나를 일대일로 대응시켜, $\ideal Q$의 어떤 거듭제곱이 $\bar{\ideal Q}$로 나누어지고 그 역도 성립하게 할 수 있다.}\footnote{서로 다른 표현의 예는 정리 II의 각주 12에 나온다. 임의의 $\mu$에 대하여 $(x^2,xy)=[(x),(x^2,\mu x+y)]$이다. 대응 소아이디얼 $\ideal P_1=(x)$, $\ideal P_2=(x,y)$가 서로 다르므로 이들은 극대 으뜸아이디얼이다. ``고립'' 극대 으뜸아이디얼의 유일성에 대해서는 \S~7을 참조하라.} \emph{또한 $\ideal Q$들과 $\bar{\ideal Q}$들은 극대 으뜸아이디얼로의 분해에 관하여 기약성 성질을 가진다.}

\textbf{따름정리.} 정리 IX는 축약표현 대신 최단표현만을 가정해도 본질적으로 그대로 성립한다는 점을 덧붙여 둔다. 예를 들어
\[
 \ideal M=[\ideal Q_1\cdots\ideal Q_i^*\cdots\ideal Q_\alpha]
\]
가 $\ideal Q_i^*$에 관하여 축약되어 있고, $\ideal Q_i^*$가 $\ideal Q_i$의 진약수이며, $\ideal L_i$가 $\ideal Q_i$의 보충부분이라고 하자. 그러면 보조정리 IV에 따라
\[
 \ideal Q_i=[\ideal Q_i^*;(\ideal L_i,\ideal Q_i)]
\]
이고, 이 표현은 $\ideal Q_i^*$에 관하여 축약되어 있다. 정리 VIII을 적용할 때 필요하면 $(\ideal L_i,\ideal Q_i)$를 진약수로 바꾸어도 되므로, $\ideal Q_i^*$는 으뜸이고 $\ideal Q_i$와 같은 대응 소아이디얼 $\ideal P_i$를 가진다. 이 절차를 계속하면 그러한 각 표현에 극대 으뜸아이디얼들로 된 축약표현을 대응시킬 수 있으며, 성분의 수와 대응 소아이디얼들은 일치한다. 따라서 최단표현의 경우에도 서로 다른 두 표현에서 성분의 수와 대응 소아이디얼들이 일치한다.

이렇게 유일하게 정의되는 서로 다른 대응 소아이디얼들을 줄여서 ``$\ideal M$에 대응 소아이디얼들''이라 부르기로 한다.

\begin{center}
\S~6.\\[0.25em]
\textbf{한 아이디얼을 상대소-기약아이디얼들의 최소공배수로 유일하게 나타내기.}
\end{center}

\textbf{정의 V.} \emph{$\ideal T\cdot\ideal R\equiv0(\ideal S)$로부터 반드시 $\ideal T\equiv0(\ideal S)$가 따르면 아이디얼 $\ideal R$을 $\ideal S$에 대하여 상대소라고 한다. $\ideal R$이 $\ideal S$에 대하여 상대소이고 $\ideal S$도 $\ideal R$에 대하여 상대소이면, $\ideal R$과 $\ideal S$를 서로 상대소라고 한다.}\footnote{상대소 조건은 대칭적이지 않다. 예를 들어 $\ideal R=(x^2,y)$는 $\ideal S=(x)$에 대하여 상대소이지만, $\ideal S$는 $\ideal R$에 대하여 상대소가 아니다. 실제로 $\ideal S^2\equiv0(\ideal R)$이지만 $\ideal T=\ideal S\not\equiv0(\ideal R)$이다.}

\emph{한 아이디얼을 서로 상대소인 두 진약수의 최소공배수로 나타낼 수 없으면 그 아이디얼을 상대소-기약이라 한다.}

특히 $\ideal T\cdot\ideal R\equiv0(\ideal S)$인 모든 $\ideal T$의 최대공약아이디얼 $\ideal T_0$를 취하면, $\ideal T_0\cdot\ideal R\equiv0(\ideal S)$이고 $\ideal S\equiv0(\ideal T_0)$이다. 따라서 $\ideal R$이 $\ideal S$에 대하여 상대소이면 $\ideal T_0=\ideal S$이고, 상대소가 아니면 $\ideal T_0$는 $\ideal S$의 진약수이다.\footnote{이렇게 정의한 $\ideal T_0$는 라스커의 ``잔여모듈''과 일치하며, 아이디얼 대신 수모듈로 확장하면 데데킨트의 두 모듈의 ``몫''과 일치한다. Lasker, 앞의 논문, 49쪽; Dedekind, \emph{Zahlentheorie}, 504쪽.}

유일성 증명의 바탕은 다음 정리이다.

\textbf{정리 X.} \emph{1. $\ideal R$이 아이디얼 $\ideal S_1,\ldots,\ideal S_\lambda$ 각각에 대하여 상대소이면, $\ideal R$은 그 최소공배수 $\ideal S$에 대해서도 상대소이다.}

\emph{2. 아이디얼 $\ideal S_1,\ldots,\ideal S_\lambda$가 각각 $\ideal R$에 대하여 상대소이면, 그 최소공배수 $\ideal S$도 $\ideal R$에 대하여 상대소이다.}

\emph{3. $\ideal R$이 $\ideal S$에 대하여 상대소이고 $\ideal S=[\ideal S_1\cdots\ideal S_\lambda]$가 $\ideal S$의 축약표현이면, $\ideal R$은 각 $\ideal S_i$에 대해서도 상대소이다.}

\emph{4. $\ideal S$가 $\ideal R$에 대하여 상대소이면, $\ideal S$의 모든 약수 $\ideal S_i$도 $\ideal R$에 대하여 상대소이다.}

1. $\ideal T\cdot\ideal R\equiv0(\ideal S)$로부터 반드시 $\ideal T\cdot\ideal R\equiv0(\ideal S_i)$가 따르므로, 가정에 따라 $\ideal T\equiv0(\ideal S_i)$이고 따라서 $\ideal T\equiv0(\ideal S)$이다.

2. 다음과 같이 놓자.
\[
 \ideal C_1=[\ideal S_2\cdots\ideal S_\lambda],\quad
 \ideal C_{12}=[\ideal S_3\cdots\ideal S_\lambda],\quad\ldots,\quad
 \ideal C_{12\ldots\lambda-1}=\ideal S_\lambda.
\]
$\ideal T\cdot\ideal S\equiv0(\ideal R)$로부터 $\ideal T\cdot\ideal S_1\cdot\ideal C_1\equiv0(\ideal R)$가 따르고, 따라서 가정에 따라 $\ideal T\cdot\ideal C_1\equiv0(\ideal R)$이다. 여기서 다시 $\ideal T\cdot\ideal S_2\cdot\ideal C_{12}\equiv0(\ideal R)$, 따라서 $\ideal T\cdot\ideal C_{12}\equiv0(\ideal R)$를 얻는다. 마침내 $\ideal T\cdot\ideal C_{12\ldots\lambda-1}=\ideal T\cdot\ideal S_\lambda\equiv0(\ideal R)$이고, 따라서 $\ideal T\equiv0(\ideal R)$이다.

3. $\ideal C_i$를 $\ideal S_i$의 보충부분이라 하자. $\ideal T_0\cdot\ideal R\equiv0(\ideal S_i)$이면 $\ideal C_i\cdot\ideal R\equiv0(\ideal S_i)$이기도 하므로 $[\ideal T_0,\ideal C_i]\cdot\ideal R\equiv0(\ideal S)$이다. 만일 $\ideal T_0$가 $\ideal S_i$의 진약수라면, 축약표현 $\ideal S=[\ideal S_i,\ideal C_i]$ 때문에 $[\ideal T_0,\ideal C_i]$도 $\ideal S$의 진약수가 되어 가정에 어긋난다.

4. $\ideal T$가 $\ideal R$의 진약수이고 $\ideal T\cdot\ideal S_i\equiv0(\ideal R)$이면 $\ideal T\cdot\ideal S\equiv0(\ideal R)$도 성립한다. 그런데 $\ideal S$가 $\ideal R$에 대하여 상대소이므로 $\ideal T\equiv0(\ideal R)$이어야 하고, 이는 $\ideal T$가 $\ideal R$의 진약수라는 가정과 모순된다.

정의 V에 따르면 특히 모든 아이디얼은 $\Sigma$의 모든 원소로 이루어진 단위아이디얼 $\ideal D$에 대하여 상대소이다.\footnote{$\ideal D$는 나눗셈과 최소공배수에 대해서만 단위의 구실을 하고, 곱셈에 대해서는 그렇지 않다. 예를 들어 상수항이 없는 $x$의 정수계수 다항식 전체의 환영역에서는 $\ideal D=(x)$이고, 모든 짝수의 환영역에서는 $\ideal D=(2)$이다.} 그러나 다음 정리들은 $\ideal D$와 다른 아이디얼들에 대해서만 성립한다.

방금 증명한 정리 X로부터 먼저 다음을 얻는다.

\textbf{보조정리 V.} \emph{한 아이디얼을 $\ideal D$와 다르며 서로 상대소인 아이디얼들의 최소공배수로 나타내는 모든 표현은 축약되어 있다.}

\[
 \ideal M=[\ideal R_1\ideal R_2\cdots\ideal R_\sigma]=[\ideal R_i,\ideal L_i]
\]
가 그러한 표현이라고 하자. 정리 X의 2에 따라 $\ideal L_i$도 $\ideal R_i$에 대하여 상대소이다. 따라서 $\ideal R_i$는 $\ideal L_i$의 약수가 될 수 없고, 이 표현은 최단이다. 실제로 $\ideal L_i\equiv0(\ideal R_i)$이면서 $\ideal L_i$가 $\ideal R_i$에 대하여 상대소라면 $\ideal D\cdot\ideal L_i\equiv0(\ideal R_i)$로부터 $\ideal D\equiv0(\ideal R_i)$, 따라서 $\ideal R_i=\ideal D$가 따르지만 이는 가정에서 배제되었다.

이제 $\ideal R_i$를 진약수 $\ideal R_i^*$로 바꿀 수 있다고 하자. 그러면 보조정리 II에 따라 $\ideal R_i$는 가약이고
\[
 \ideal R_i=[\ideal R_i^*;(\ideal R_i,\ideal L_i)]
\]
이다. 여기서 필요하면 $(\ideal R_i,\ideal L_i)$를 진약수 $(\ideal R_i,\ideal L_i)^*$로 바꾸고 $\ideal R_i^*$도 진약수로 바꾸어 $\ideal R_i$의 축약표현을 얻는다. 그러면 정리 X의 3에 따라 $\ideal L_i$는 $(\ideal R_i,\ideal L_i)^*$에 대해서도 상대소이고, 이 아이디얼은 최단표현 때문에 $\ideal D$와 다르다. 그런데 $(\ideal R_i,\ideal L_i)^*$는 $(\ideal R_i,\ideal L_i)$의 약수이고 $(\ideal R_i,\ideal L_i)$는 $\ideal L_i$의 약수이므로 모순이다.

정리 X로부터 대응 소아이디얼과의 관계를 주는 다음 정리도 얻는다.

\textbf{정리 XI.} \emph{$\ideal R$이 $\ideal S$에 대하여 상대소이고 $\ideal S$가 $\ideal D$와 다르면, $\ideal R$에 결부된 어느 소아이디얼도\footnote{$\ideal R$에 대응 소아이디얼의 정의는 \S~5의 끝을 보라.} $\ideal S$에 대응 소아이디얼로 나누어지지 않는다. 거꾸로 그러한 나눗셈 관계가 하나도 없으면 $\ideal R$은 $\ideal S$에 대하여 상대소이고, 물론 $\ideal S$는 $\ideal D$와 다르다.}

\[
 \ideal R=[\ideal Q_1\cdots\ideal Q_\alpha],\qquad
 \ideal S=[\ideal Q_1^*\cdots\ideal Q_\beta^*]
\]
가 $\ideal R$과 $\ideal S$를 극대 으뜸아이디얼들로 나타내는 축약표현이고, $\ideal P_1,\ldots,\ideal P_\alpha$와 $\ideal P_1^*,\ldots,\ideal P_\beta^*$가 각각 대응 소아이디얼들이라고 하자. 우리는 다음과 같은 형태로 주장을 보인다. 어떤 $\ideal P$가 어떤 $\ideal P^*$로 나누어지면 $\ideal R$은 $\ideal S$에 대하여 상대소일 수 없으며, 그 역도 성립한다.

실제로 $\ideal P_\mu\equiv0(\ideal P_\nu^*)$라고 하자. 그러면 $\ideal Q_\mu\equiv0(\ideal P_\nu^*)$이기도 하고, $\ideal P_\nu^*$의 정의에 따라 $\ideal Q_\mu^{\sigma_\nu}\equiv0(\ideal Q_\nu^*)$, 따라서 $\ideal R^{\sigma_\nu}\equiv0(\ideal Q_\nu^*)$이다. $\ideal R$의 거듭제곱 가운데 $\ideal Q_\nu^*$로 나누어지는 가장 낮은 것을 $\ideal R^\tau$라고 하자. $\tau=1$이면 $\ideal R$이 $\ideal Q_\nu^*$로 나누어진다. 그런데 가정에 따라 $\ideal S$가 $\ideal D$와 다르므로 $\ideal R$은 $\ideal Q_\nu^*$에 대하여 상대소가 아니다. $\tau\ge2$이면
\[
 \ideal R^{\tau-1}\cdot\ideal R\equiv0(\ideal Q_\nu^*),\qquad
 \ideal R^{\tau-1}\not\equiv0(\ideal Q_\nu^*)
\]
이므로 역시 $\ideal R$은 $\ideal Q_\nu^*$에 대하여 상대소가 아니다.\footnote{$\Sigma$가 단위원을 포함할 필요가 없으므로 $\ideal R^0$은 정의되지 않는다. 따라서 $\tau=1$인 경우를 먼저 따로 다루어야 했다. $\Sigma$가 단위원을 포함하는 경우에도 $\tau=0$은 $\ideal S$에 관한 가정 때문에 배제된다.} 그러므로 두 경우 모두 정리 X의 3에 따라 $\ideal R$은 $\ideal S$에 대해서도 상대소가 아니다.

거꾸로 $\ideal R$이 $\ideal S$에 대하여 상대소가 아니면, 정리 X의 1에 따라 $\ideal R$은 적어도 하나의 $\ideal Q_\nu^*$에 대하여 상대소가 아니다. 따라서
\[
 \ideal T_0\cdot\ideal R\equiv0(\ideal Q_\nu^*),\qquad
 \ideal T_0\not\equiv0(\ideal Q_\nu^*)
\]
이고, $\ideal Q_\nu^*$가 으뜸이므로
\[
 \ideal R^\tau\equiv0(\ideal Q_\nu^*),\qquad
 \text{따라서}\qquad
 \ideal Q_1^\tau\cdots\ideal Q_\alpha^\tau\equiv0(\ideal Q_\nu^*)
\]
이다. 그러면 대응 소아이디얼들에 대하여
\[
 \ideal P_1^{\tau\varrho_1}\cdots\ideal P_\alpha^{\tau\varrho_\alpha}
 \equiv0(\ideal P_\nu^*)
\]
를 얻는다. 소아이디얼의 성질에 따라 적어도 하나의 $\ideal P$가 $\ideal P_\nu^*$로 나누어진다. 이로써 정리 XI가 증명되었다.

이제 정리 X와 XI로부터 상대소-기약아이디얼로의 분해가 존재하고\footnote{분해의 존재는 \S~2에서 유한개의 기약아이디얼로의 분해가 존재함을 증명한 것과 정확히 같은 방식으로 직접 증명할 수도 있다.} 유일함을 다음과 같이 얻는다.

$\ideal M=[\ideal Q_1\cdots\ideal Q_\alpha]$가 $\ideal M$을 극대 으뜸아이디얼들로 나타내는 축약표현(또는 적어도 최단표현)이고, $\ideal P_1,\ldots,\ideal P_\alpha$가 대응 소아이디얼들이라고 하자. 어느 한 군의 아이디얼도 다른 군의 아이디얼로 나누어지지 않도록 $\ideal P$들을 군들로 묶되, 각 군 자체는 이 성질을 갖는 두 부분군으로 더 나눌 수 없게 한다.

그러한 군 분할을 구성하려면, 정의상 각 군 $G$가 한 아이디얼 $\ideal P$와 함께 $\ideal P_1,\ldots,\ideal P_\alpha$ 가운데 나타나는 그 모든 약수와 배수(즉 $\ideal P$로 나누어지는 아이디얼)를 포함해야 함을 관찰하면 된다. 이를테면 $\ideal P^{(i_1)}$들을 $\ideal P$의 모든 배수, $\ideal P_{j_1}^{(i_1)}$들을 $\ideal P^{(i_1)}$의 모든 약수, $\ideal P_{j_1}^{(i_1i_2)}$들을 $\ideal P_{j_1}^{(i_1)}$의 모든 배수 등으로 정한다. 일반적으로 $\ideal P_{j_1\ldots j_{\lambda-1}}^{(i_1\ldots i_\lambda)}$들을 $\ideal P_{j_1\ldots j_{\lambda-1}}^{(i_1\ldots i_{\lambda-1})}$의 모든 배수로, $\ideal P_{j_1\ldots j_\lambda}^{(i_1\ldots i_\lambda)}$들을 $\ideal P_{j_1\ldots j_{\lambda-1}}^{(i_1\ldots i_\lambda)}$의 모든 약수로 정한다. 전체적으로 유한개의 아이디얼 $\ideal P$만을 다루므로 이 절차는 유한 단계 뒤에 끝나야 한다. 곧 앞서 나온 모든 아이디얼과 다른 새 아이디얼을 더 이상 내놓지 않는다. 이렇게 얻은 아이디얼 $\ideal P$들의 전체는 실제로 원하는 성질을 가진 한 군 $G$를 이룬다.

실제로 정의에 따라 $G$는 각 $\ideal P_{j_1\ldots j_{\lambda-1}}^{(i_1\ldots i_{\lambda-1})}$와 함께 그 모든 배수 $\ideal P_{j_1\ldots j_{\lambda-1}}^{(i_1\ldots i_\lambda)}$를 포함하고, 각 $\ideal P_{j_1\ldots j_{\lambda-1}}^{(i_1\ldots i_\lambda)}$와 함께 그 모든 약수 $\ideal P_{j_1\ldots j_\lambda}^{(i_1\ldots i_\lambda)}$를 포함한다. 여기에는 $\ideal P_{j_1\ldots j_{\lambda-1}}^{(i_1\ldots i_{\lambda-1})}$의 모든 약수도 들어 있고, $\ideal P_{j_1\ldots j_{\lambda-1}}^{(i_1\ldots i_\lambda)}$의 배수들도 다시 같은 방식으로 들어 있다. 따라서 $G$에 들어 있지 않은 아이디얼은 $G$에 들어 있는 아이디얼의 약수도 배수도 될 수 없다.

$G$는 기약성 조건도 만족한다. 실제로 $G$가 두 부분군 $G^{(1)}$, $G^{(2)}$로 나뉘고, 이를테면 $G^{(1)}$이 $\ideal P_{j_1\ldots j_\lambda}^{(i_1\ldots i_\lambda)}$(또는 $\ideal P_{j_1\ldots j_{\lambda-1}}^{(i_1\ldots i_\lambda)}$)를 포함한다고 하자. 그러면 앞선 아이디얼들은 번갈아 그 배수와 약수(또는 약수와 배수)이므로 $G^{(1)}$은 그 모두를 포함하고, 따라서 $\ideal P$와 군 $G$ 전체를 포함한다. $G$에 들어 있지 않은 아이디얼들에 같은 절차를 적용하면 모든 $\ideal P$의 군 분할 $G_1,\ldots,G_\sigma$를 얻으며, 이 분할은 원하는 성질들을 가진다. 그러한 군 분할은 유일하다. 실제로 $G'_1,\ldots,G'_\tau$가 두 번째 분할이고 $\ideal P_{j_1\ldots j_\lambda}^{(i_1\ldots i_\lambda)}$(또는 $\ideal P_{j_1\ldots j_{\lambda-1}}^{(i_1\ldots i_\lambda)}$)가 $G'_i$의 원소이면, 위의 논증에 따라 $G'_i$는 군 $G$ 전체를 포함하고, 기약성 성질에 따라 $G$와 일치한다.

$\ideal Q_{i\mu}$($\mu=1,\ldots,\lambda_i$)가 군 $G_i$에 함께 묶인 아이디얼 $\ideal Q$들을 나타낸다고 하자. $\ideal P$들이 모두 서로 다르므로 최단표현에서 이들에 대응하는 으뜸아이디얼 $\ideal Q_{i\mu}$들은 유일하게 정해진다. 이제
\[
 \ideal R_i=[\ideal Q_{i1}\cdots\ideal Q_{i\lambda_i}],\qquad
 \text{그러면}\qquad
 \ideal M=[\ideal R_1\cdots\ideal R_\sigma]
\]
로 놓자. 이것이 $\ideal M$의 상대소-기약아이디얼 분해임을 보이겠다. 먼저 $\ideal R_i=[\ideal Q_{i1}\cdots\ideal Q_{i\lambda_i}]$가 최단표현일 뿐이어도 정리 XI를 적용할 수 있음을 관찰하자. 정리 IX의 따름정리에 따라 대응 소아이디얼들이 이미 유일하게 정의되기 때문이다. 이제 $\ideal R_i$에 결부된 어느 소아이디얼 $\ideal P_{i\mu}$도 $\ideal R_j$에 대응 소아이디얼 $\ideal P_{j\nu}$로 나누어지지 않고, 그 역도 성립하므로 정리 XI에 따라 $\ideal R_i$와 $\ideal R_j$는 서로 상대소이다. 또한 각 $\ideal R_i$는 군 $G_i$의 기약성 성질 때문에 정리 XI에 따라 상대소-기약이다. 가약인 경우에는 언제나 $\ideal D$와 다른 아이디얼들만을 다루기 때문이다. 더 나아가 처음에 $\ideal Q$들의 최단표현만에서 출발했더라도 보조정리 V에 따라 이 표현은 축약되어 있다. 거꾸로 정리 XI에 따라 상대소-기약아이디얼로의 모든 분해는 위에서 설명한 $\ideal P_{i\mu}$들의 군 분할을 낳는다.

이제
\[
 \ideal M=[\bar{\ideal R}_1\cdots\bar{\ideal R}_\tau]
\]
가 $\ideal M$을 상대소-기약아이디얼들로 나타내는 두 번째 표현이라고 하자. 이 표현은 보조정리 V에 따라 축약되어 있다. $\ideal R$들을 극대 으뜸아이디얼들로 분해해 보면 대응 소아이디얼들이 일치하므로, 위에서 유일성을 증명한 이 소아이디얼들의 군 분할도 일치한다. 따라서 $\tau=\sigma$이고, $\ideal R_i$와 $\bar{\ideal R}_i$가 같은 군에 속하도록 지표를 정할 수 있다. 이제 아이디얼과 보충부분으로 된 표현을
\[
 \ideal M=[\ideal R_i,\ideal L_i]
          =[\bar{\ideal R}_i,\bar{\ideal L}_i]
\]
라 하자. $\bar{\ideal R}_i$가 $\ideal R_i$와 같은 군에 대응하므로 정리 XI에 따라 $\ideal L_i$는 $\bar{\ideal R}_i$에 대하여 상대소이고, $\bar{\ideal L}_i$는 $\ideal R_i$에 대하여 상대소이다. 또한
\[
 \ideal R_i\cdot\ideal L_i\equiv0(\bar{\ideal R}_i),
 \qquad
 \bar{\ideal R}_i\cdot\bar{\ideal L}_i\equiv0(\ideal R_i)
\]
이므로
\[
 \ideal R_i\equiv0(\bar{\ideal R}_i),
 \qquad
 \bar{\ideal R}_i\equiv0(\ideal R_i),
 \qquad
 \ideal R_i=\bar{\ideal R}_i
\]
를 얻는다. 이로써 다음이 증명되었다.

\textbf{정리 XII.} \emph{각 아이디얼은 유한개의 서로 상대소인 상대소-기약아이디얼들의 최소공배수로 유일하게 나타낼 수 있다.}

\begin{center}
\S~7.\\[0.25em]
\textbf{고립아이디얼의 유일성.}
\end{center}

\textbf{정의 VI.} \emph{최단표현 $\ideal M=[\ideal R,\ideal L]$가 $\ideal L$에 관하여 축약되어 있다고 하자. $\ideal R$에 결부된 어느 소아이디얼도 $\ideal L$에 대응 소아이디얼의 약수가 아니면, 다시 말해 $\ideal L$이 $\ideal R$에 대하여 상대소이면 $\ideal R$을 고립아이디얼이라 한다.}

그러면 표현 $\ideal M=[\ideal R,\ideal L]$은 보조정리 V의 표현 $\ideal M=[\ideal R_i,\ideal L_i]$에 대한 조건을 만족하므로, 보조정리 V와 함께 다음이 증명된다.

\textbf{보조정리 VI.} \emph{$\ideal R$이 $\ideal L$에 관하여 축약된 최단표현 $\ideal M=[\ideal R,\ideal L]$의 고립아이디얼이면, 이 표현은 $\ideal R$에 관해서도 축약되어 있다.}

따라서 고립아이디얼의 경우에는 축약표현을 다루므로, $\ideal R$과 $\ideal L$을 기약아이디얼들로 분해할 때 나타나는 대응 소아이디얼들을 합치면 그에 대응하는 $\ideal M$의 분해에서 나타나는 유일하게 정해진 대응 소아이디얼들이 된다. 그러므로 $\ideal R$의 기약아이디얼 분해에 속하는 어느 소아이디얼도, 그에 대응하는 $\ideal M$의 분해에 속하는 나머지 소아이디얼의 약수가 아니다. 거꾸로 이 조건이 성립하고 $\ideal R$이 $\ideal R$에 관하여 축약된 표현 $\ideal M=[\ideal R,\ideal L]$에 적어도 한 번 나타나면, 따라서 어떤 축약표현 $\ideal M=[\ideal R,\ideal L^*]$에도 나타나면, 정의 VI에 따라 $\ideal R$은 고립되어 있다. 이로부터 특정 보충부분 $\ideal L$에 의존하지 않는 다음 정의를 얻는다.

\textbf{정의 VIa.} \emph{$\ideal R$의 기약아이디얼 분해에 속하는 소아이디얼들이 그에 대응하는 $\ideal M$의 분해에 속하는 나머지 소아이디얼들의 약수가 아니고, $\ideal R$이 $\ideal R$에 관하여 축약된 표현 $\ideal M=[\ideal R,\ideal L]$에 적어도 한 번 나타나면 $\ideal R$을 고립아이디얼이라 한다.}\footnote{통상적인 대응 소아이디얼들(\S~5 끝)을 사용한다면, $\ideal L$에 대응 소아이디얼들이 모두 $\ideal R$에 대응 소아이디얼들과 다르다는 조건을 따로 덧붙여야 한다. 따라서 정의 VIa에서는 표현이 보충부분에 관하여 축약되어 있다고 더 이상 가정할 필요가 없다.}

이제
\[
 \ideal M=[\ideal R,\ideal L]=[\bar{\ideal R},\bar{\ideal L}]
\]
가 고립아이디얼 $\ideal R,\bar{\ideal R}$와 보충부분 $\ideal L,\bar{\ideal L}$로 된 $\ideal M$의 두 표현이고, $\ideal R$과 $\bar{\ideal R}$에 대응 소아이디얼들이 일치한다고 하자. 표현들이 축약되도록 $\ideal L$과 $\bar{\ideal L}$을 각각 약수 $\ideal L^*$와 $\bar{\ideal L}^{*}$로 바꾸면, $\ideal L^*$와 $\bar{\ideal L}^{*}$에 대응 소아이디얼들도 일치한다. 따라서 정리 XI에 따라 $\ideal L^*$는 $\bar{\ideal R}$에 대하여 상대소이고 $\bar{\ideal L}^{*}$는 $\ideal R$에 대하여 상대소이다. 또한
\[
 \ideal R\cdot\ideal L^*\equiv0(\bar{\ideal R}),
 \qquad
 \bar{\ideal R}\cdot\bar{\ideal L}^{*}\equiv0(\ideal R)
\]
이므로
\[
 \ideal R\equiv0(\bar{\ideal R}),
 \qquad
 \bar{\ideal R}\equiv0(\ideal R),
 \qquad
 \ideal R=\bar{\ideal R}
\]
를 얻는다. 따라서 고립아이디얼들은 대응 소아이디얼들에 의해 유일하게 정해진다. 특히 이것은 기약아이디얼 또는 극대 으뜸아이디얼로의 분해에 관한 정리 VII과 IX를 강화한다. 정리 IX의 따름정리에 따라 여기서는 최단표현만을 가정해도 된다. 정리하면 다음과 같다.

\textbf{정리 XIII.} \emph{한 아이디얼을 기약아이디얼들 또는 극대 으뜸아이디얼들의 최소공배수로 나타내는 모든 최단표현에서, 고립된 기약아이디얼들 또는 고립된 극대 으뜸아이디얼들은 유일하게 정해진다. 다의성은 고립되지 않은 기약아이디얼들 또는 극대 으뜸아이디얼들에만 관계한다.}\footnote{다항식의 아이디얼을 극대 으뜸아이디얼들로 분해하는 경우에 이 정리는 이미 매콜리가 증명 없이 제시하였다. 고립 및 비고립(매입된) 으뜸아이디얼에 대한 그의 정의는 아래에서 제시하는 것의 비유리적 정식화로 볼 수 있다.} \emph{일반적으로 고립아이디얼들은 대응 소아이디얼들에 의해 유일하게 정해진다.}

그러한 기약아이디얼 또는 극대 으뜸아이디얼의 최단표현에서 아이디얼 $\ideal B_i$ 또는 $\ideal Q_j$가 고립되어 있지 않다고 하자. 그러면 정의에 따라 보충부분 $\ideal U_i$ 또는 $\ideal L_j$는 $\ideal P_i$ 또는 $\ideal P_j$로 나누어지고, 따라서
\[
 \ideal U_i^{\varrho_i}\equiv0(\ideal B_i)
 \qquad\text{또는}\qquad
 \ideal L_j^{\sigma_j}\equiv0(\ideal Q_j)
\]
이다. 거꾸로 이 관계들이 성립하면 $\ideal P_i$ 또는 $\ideal P_j$는 보충부분에 결부된 적어도 하나의 소아이디얼의 약수이다. 따라서 정리 IX의 따름정리에 따라, $\ideal L_j$의 약수로서 $\ideal L_j^*$에 관하여 축약된 표현을 주는 $\ideal L_j^*$에 대응 소아이디얼 하나의 약수이기도 하다. 그러므로 $\ideal B_i$ 또는 $\ideal Q_j$는 고립되어 있지 않다. 특히 $\ideal M$의 분해에서 대응 소아이디얼이 한 번보다 많이 나타나는 기약아이디얼 $\ideal B_i$는 언제나 고립되어 있지 않다. \emph{따라서 고립되지 않은 으뜸아이디얼은 모든 보충부분의 어떤 거듭제곱이 그 아이디얼로 나누어진다는 성질로도 특징지어지고, 고립된 으뜸아이디얼은 이 성질이 성립할 수 없다는 것으로 특징지어진다.}

\begin{center}
\S~8.\\[0.25em]
\textbf{한 아이디얼을 서로소-기약아이디얼들의 곱으로 유일하게 나타내기.}
\end{center}

바탕 환영역 $\Sigma$가 단위원, 즉 $\Sigma$의 모든 원소 $a$에 대하여 $\varepsilon\cdot a=a$가 되는 원소 $\varepsilon$을 가진다고 하자.\footnote{곱셈이 가환이므로 $\Sigma$는 단위원을 둘 이상 가질 수 없다. 실제로 두 단위원 $\varepsilon_1,\varepsilon_2$에 대하여 $\varepsilon_1\varepsilon_2=\varepsilon_2=\varepsilon_1$이다.} 그러면 서로소 아이디얼들을 다음과 같이 정의할 수 있다.

\textbf{정의 VIII.} \emph{두 아이디얼 $\ideal R,\ideal S$의 최대공약아이디얼이 $\Sigma$의 모든 원소로 이루어진 단위아이디얼 $\ideal D=(\varepsilon)$와 같으면 $\ideal R$과 $\ideal S$를 서로소라고 한다. 한 아이디얼을 쌍마다 서로소인 아이디얼들의 최소공배수로 나타낼 수 없으면 그 아이디얼을 서로소-기약이라 한다.}

서로소인 두 아이디얼은 언제나 서로 상대소임을 관찰하자. 실제로 정의에 따라 $r\equiv0(\ideal R)$, $s\equiv0(\ideal S)$이고 $\varepsilon=r+s$인 두 원소 $r,s$가 존재한다. $\ideal T\cdot\ideal R\equiv0(\ideal S)$이면 $\ideal T\cdot r\equiv0(\ideal S)$이고, 따라서 $\ideal T\cdot\varepsilon=\ideal T\equiv0(\ideal S)$이다. 마찬가지로 $\bar{\ideal T}\cdot\ideal S\equiv0(\ideal R)$이면 $\bar{\ideal T}\equiv0(\ideal R)$이다.\footnote{그러나 역은 성립하지 않는다. 예를 들어 아이디얼 $\ideal R=(x)$와 $\ideal S=(y)$는 서로 상대소이지만 서로소는 아니다.} 따라서 보조정리 V에 따라 쌍마다 서로소인 아이디얼들로 된 모든 표현은 동시에 축약되어 있다.

유일성 증명의 바탕은 정리 X에 대응하는 다음 정리이다.

\textbf{정리 XIV.} \emph{$\ideal R$이 아이디얼 $\ideal S_1,\ldots,\ideal S_\lambda$ 각각과 서로소이면, $\ideal R$은 $\ideal S=[\ideal S_1\cdots\ideal S_\lambda]$와도 서로소이다. 거꾸로 $\ideal R$과 $\ideal S$가 서로소이면 $\ideal R$은 각 $\ideal S_j$와도 서로소이다. $\ideal R=[\ideal R_1\cdots\ideal R_\mu]$이고 각 $\ideal R_i$가 각 $\ideal S_j$와 서로소이면 $\ideal R$과 $\ideal S$는 서로소이며, 여기서도 역이 성립한다.}

실제로 $\ideal R$이 각 $\ideal S_j$와 서로소이면
\[
 s_j\equiv0(\ideal S_j),
 \qquad
 s_j\equiv\varepsilon(\ideal R)
\]
인 원소 $s_j$들이 존재한다. 따라서
\[
 s_1\cdot s_2\cdots s_\lambda\equiv0(\ideal S),
 \qquad
 s_1\cdot s_2\cdots s_\lambda\equiv\varepsilon(\ideal R),
 \qquad
 (\ideal R,\ideal S)=(\varepsilon)
\]
이다. 한편 $(\ideal R,\ideal S)$는 각 $(\ideal R,\ideal S_j)$로 나누어지므로 역도 성립한다. 이 논증을 여러 번 적용하면 주장의 둘째 부분을 얻는다. 실제로 고정된 $i$에 대하여 $\ideal R_i$가 $\ideal S_1,\ldots,\ideal S_\lambda$ 각각과 서로소이면 $\ideal R_i$는 $\ideal S$와 서로소이다. 이 사실이 모든 $i$에 대하여 성립하면 서로소 개념이 상호적이므로 $\ideal S$는 $\ideal R$과 서로소이다. 거꾸로 $\ideal R$과 $\ideal S$가 서로소이면 $\ideal S$는 각 $\ideal R_i$와 서로소이고, 따라서 $\ideal R_i$는 $\ideal S$와 서로소이다.

서로소-기약아이디얼로의 분해가 존재하고 유일함을 보이는 증명은\footnote{분해의 존재는 다시 \S~2와 비슷하게 직접 증명할 수 있고, 유일성도 직접 증명할 수 있다(서론에서 슈마이들러와 뇌터--슈마이들러에 관하여 말한 내용을 참조하라). 여기서 제시하는 증명은 서로소-기약아이디얼의 구조도 함께 밝혀 준다.} 상대소아이디얼의 경우에 대응하는 증명처럼 유일한 군 분할로 귀착된다. 그러나 $\ideal M$의 상대소-기약아이디얼 $\ideal R_1,\ldots,\ideal R_\sigma$는 정리 XII에 따라 유일하게 정해져 있으므로, 여기서는 대응 소아이디얼들로 거슬러 올라갈 필요가 없다.

$\ideal M$의 유일하게 정해진 상대소-기약아이디얼 $\ideal R_1,\ldots,\ideal R_\sigma$를 다음과 같이 군들로 묶는다. \emph{한 군의 각 아이디얼이 다른 군의 각 아이디얼과 서로소가 되게 하되, 각 군 자체는 한 부분군의 각 아이디얼이 다른 부분군의 각 아이디얼과 서로소가 되도록 두 부분군으로 나눌 수 없게 한다.} 그러한 군 분할은 다음과 같이 얻는다. 정의상 각 군 $G$는 한 아이디얼 $\ideal R$과 함께 $\ideal R$과 서로소가 아닌 모든 아이디얼도 포함해야 한다. 이들을 $\ideal R_{i_1}$이라 하자. 다시 이들과 서로소가 아닌 아이디얼들을 $\ideal R_{i_1i_2}$라 하고, 일반적으로 $\ideal R_{i_1\ldots i_\lambda}$가 $\ideal R_{i_1\ldots i_{\lambda-1}}$과 서로소가 아니라고 하자. 전체적으로 유한개의 아이디얼만을 다루므로 이 절차는 유한 단계 뒤에 끝나야 한다. 곧 앞서 나온 모든 아이디얼과 다른 새 아이디얼을 더 이상 내놓지 않는다. 이렇게 얻은 아이디얼들의 전체는 원하는 성질을 가진 한 군 $G$를 이룬다. 실제로 $G$에 들어 있지 않은 모든 아이디얼은 구성에 따라 $G$에 들어 있는 모든 아이디얼과 서로소이다. 더 나아가 $G$가 두 부분군 $G^{(1)}$, $G^{(2)}$로 나뉘고 이를테면 $\ideal R_{i_1\ldots i_\lambda}$가 $G^{(1)}$의 원소라고 하자. 서로소성은 상호적이므로 $G^{(1)}$은 $\ideal R_{i_1\ldots i_{\lambda-1}},\ldots,\ideal R_{i_1}$과 $\ideal R$, 따라서 군 $G$ 전체도 포함해야 한다. 이로써 기약성이 증명되었다. $G$에 들어 있지 않은 아이디얼들에 같은 절차를 적용하면 모든 $\ideal R$의 군 분할 $G_1,\ldots,G_\tau$를 얻는다.

이 군 분할은 유일하다. 실제로 $G'_1,\ldots,G'_\tau$가 두 번째 분할이고 $\ideal R_{i_1\ldots i_\lambda}$가 $G'_i$의 원소라고 하자. 그러면 $G'_i$는 $\ideal R$과 따라서 $G_1$도 포함하며, 기약성 조건 때문에 $G_1$과 다른 원소를 포함할 수 없다. 그러므로 $G'_i=G_1$이다.

이제 $\ideal T_i$를 군 $G_i$에 모인 아이디얼 $\ideal R$들의 최소공배수라 하자. 그러면
\[
 \ideal M=[\ideal T_1\cdots\ideal T_\tau]
\]
가 $\ideal M$을 서로소-기약아이디얼들로 나타내는 표현임을 보이겠다. 먼저 $\ideal T$들을 $\ideal R$들로 분해해 보면 $[\ideal T_1\cdots\ideal T_\tau]$가 실제로 $\ideal M$을 나타냄을 알 수 있다. 또한 정리 XIV에 따라 $\ideal T$들은 쌍마다 서로소이고, 각 $\ideal T$는 서로소-기약이다. 이로써 그러한 표현의 존재가 증명되었다.

유일성을 증명하기 위하여 $\ideal M=[\bar{\ideal T}_1\cdots\bar{\ideal T}_\tau]$가 그러한 두 번째 표현이라고 하자. $\bar{\ideal T}$들을 상대소-기약아이디얼 $\ideal R$들로 분해하면 서로 다른 $\bar{\ideal T}$에서 나타나는 $\ideal R$들은 정리 XIV에 따라 서로소이고, 따라서 서로 상대소이기도 하다. 그러므로 이들은 $\ideal M$의 유일하게 정해진 상대소-기약아이디얼 $\ideal R$들과 일치한다. 또한 정리 XIV에 따라 $\bar{\ideal T}_i$들은 앞서 말한 성질을 가진 $\ideal R$들의 군 분할 $G'_i$를 만든다. 그런데 이 군 분할은 유일하고 각 $\bar{\ideal T}_i$는 군 $G'_i=G_i$에 의해 유일하게 정해지므로 $\bar{\ideal T}_i=\ideal T_i$이다. 이로써 유일성이 증명되었다.

쌍마다 서로소인 아이디얼들에서는 최소공배수가 곱과 같다. 실제로 정리 XIV에 따라 $\ideal M=[\ideal T_1\cdots\ideal T_\tau]$에서 보충부분 $\ideal L_i$도 $\ideal T_i$와 서로소이다. 따라서
\[
 t_i\equiv0(\ideal T_i),
 \qquad
 l_i\equiv0(\ideal L_i),
 \qquad
 \varepsilon=t_i+l_i
\]
인 원소들이 존재한다. $f\equiv0(\ideal T_i)$, $f\equiv0(\ideal L_i)$이면
\[
 f=f\varepsilon=ft_i+fl_i
 \quad\text{이므로}\quad
 f\equiv0(\ideal L_i\cdot\ideal T_i)
\]
이다. 거꾸로 $\ideal L_i\cdot\ideal T_i$는 $[\ideal L_i,\ideal T_i]$로 나누어지므로 $[\ideal L_i,\ideal T_i]=\ideal L_i\cdot\ideal T_i$이다. 이 절차를 $\ideal L_i$에 계속 적용하면 마침내
\[
 \ideal M=[\ideal T_1\cdots\ideal T_\tau]
       =\ideal T_1\cdot\ideal T_2\cdots\ideal T_\tau
\]
를 얻는다. 이로써 다음이 증명되었다.

\textbf{정리 XV.} \emph{각 아이디얼은 유한개의 쌍마다 서로소인 서로소-기약아이디얼들의 곱으로 유일하게 나타낼 수 있다.}

\begin{center}
\S~9.\\[0.25em]
\textbf{모듈로의 논의 확장. 기약모듈 분해에서 성분 수의 동일성.}
\end{center}

이제 으뜸아이디얼이나 소아이디얼이 아니라 \emph{기약아이디얼}에 관한 앞의 세 절의 내용이 더 약한 가정 아래에서도 성립함을 보이겠다. 실제로 그 절들은 곱셈의 가환법칙을 사용하지 않고 아이디얼이 모듈이라는 성질만을 사용한다. 따라서 이제 정의할 비가환 환영역에 관한 모듈에도 그대로 성립한다. 모듈을 정의하려면 다음 성질들을 가진 \emph{이중환영역 $(\Sigma,T)$}를 바탕으로 삼아야 한다.

\emph{$\Sigma$는 추상적으로 정의된 비가환환이다.} 즉 $\Sigma$는 원소 $a,b,c,\ldots$의 체계이고, 이 원소들에는 환의 덧셈 $(\#)$과 환의 곱셈 $(\times)$이라는 두 연산이 정의되어 있다. 이 연산들은 \S~1에서 세운 법칙들 가운데 환의 곱셈에 관한 가환법칙 4를 제외한 모든 법칙을 만족한다.

$T$는 원소 $\alpha,\beta,\gamma,\ldots$의 체계이며, $\Sigma$와 결합하여 여기에도 두 연산이 정의되어 있다. \emph{덧셈}은 두 원소 $\alpha,\beta$에서 세 번째 원소 $\alpha+\beta$를 유일하게 만들고, \emph{$T$의 원소 $\alpha$와 $\Sigma$의 원소 $c$의 곱셈}은 $T$의 원소 $c\cdot\alpha$를 유일하게 만든다.\footnote{따라서 여기서는 ``우측'' 곱셈, ``우측'' 환영역 $T$, 그리고 그에 따라 ``우측'' 모듈과 아이디얼을 다룬다. $T$에 좌측 곱셈 $\alpha\cdot c$를 바탕으로 삼으면 좌측모듈과 좌측아이디얼의 대응하는 이론을 얻는다. 즉 $M$은 $\alpha$와 함께 $\alpha\cdot c$도 포함한다. 이때 결합법칙은 $(\gamma\cdot b)\cdot a=\gamma\cdot(b\times a)$ 꼴이다.}

이 연산들에는 다음 법칙들이 성립한다.
\begin{enumerate}
\item 덧셈의 결합법칙: $(\alpha+\beta)+\gamma=\alpha+(\beta+\gamma)$.
\item 덧셈의 가환법칙: $\alpha+\beta=\beta+\alpha$.
\item 제한 없고 유일한 뺄셈의 법칙: 방정식 $\alpha+\xi=\beta$를 만족하는 원소 $\xi$가 $T$에 유일하게 존재한다($\xi=\beta-\alpha$로 쓴다).
\item 곱셈의 결합법칙: $a\cdot(b\cdot\gamma)=(a\times b)\cdot\gamma$.
\item 분배법칙:
\[
 (a\# b)\cdot\gamma=a\cdot\gamma+b\cdot\gamma,
 \qquad c\cdot(\alpha+\beta)=c\cdot\alpha+c\cdot\beta.
\]
\end{enumerate}
이 조건들로부터 영원소의 존재와 뺄셈과 곱셈의 결합에 대해서도 분배법칙
\[
 (a-b)\cdot\gamma=a\cdot\gamma-b\cdot\gamma,
 \qquad c\cdot(\alpha-\beta)=c\cdot\alpha-c\cdot\beta
\]
이 성립함은 잘 알려져 있다. 여기서 $(-)$는 $\Sigma$에서의 뺄셈이다. \emph{$\Sigma$가 단위원 $\varepsilon$을 포함하면 $T$의 모든 원소 $\alpha$에 대하여 $\varepsilon\cdot\alpha=\alpha$라고 한다.}

$(\Sigma,T)$의 모듈 $M$은 다음 두 조건을 만족하는 $T$의 원소들의 체계를 뜻한다.
\begin{enumerate}
\item \emph{$M$은 $\alpha$와 함께 $c\cdot\alpha$도 포함한다. 여기서 $c$는 $\Sigma$의 임의의 원소이다.}
\item \emph{$M$은 $\alpha,\beta$와 함께 차 $\alpha-\beta$도 포함한다.} 따라서 $\alpha$와 함께 모든 정수 $n$에 대한 $n\alpha$도 포함한다.\footnote{정수는 다시 환의 원소가 아니라 줄여 쓰는 기호로 이해해야 한다.}
\end{enumerate}
이 정의에 따라 \emph{$T$ 자체도 $(\Sigma,T)$의 모듈을 이룬다.} 특히 환영역 $T$와 거기서 정한 연산들이 환영역 $\Sigma$ 및 그 연산들과 일치하면 모듈 $M$은 $\Sigma$의 우측아이디얼 $M$이 된다. 더 나아가 $\Sigma$가 가환이라고 가정하면 통상적인 아이디얼 개념을 얻는다. 따라서 아이디얼은 모듈 개념의 특수한 경우이다.\footnote{가장 간단한 모듈의 예는 정수계수 선형형식들의 모듈이다. 여기서 $\Sigma$는 모든 유리정수로, $T$는 모든 정수계수 선형형식으로 이루어진다. $\Sigma$와 $T$에서 유리정수 대신 대수적 정수를 쓰거나, 이를테면 모든 짝수를 쓰면 조금 더 일반적인 모듈을 얻는다. 선형형식 대신 매번 계수 전체를 하나의 원소로 보면 $\Sigma$와 $T$의 연산들은 실제로 서로 다르다. 비가환 다항식 환영역의 아이디얼은 뇌터--슈마이들러의 공동논문에서 다룬다. 다른 특수한 비가환 환영역의 아이디얼은 후르비츠의 \emph{사원수의 수론 강의}(베를린, 슈프링어, 1919)와 거기서 인용한 뒤 파스키에의 논문들에서 다룬다.}

모듈에 대해서도 \S~1의 모든 정의가 유지된다. 따라서 $\alpha\eqzero{M}$ 또는 $N\eqzero{M}$은 $\alpha$ 또는 $N$의 모든 원소가 $M$의 원소임을 뜻한다. 달리 말해 $\alpha$ 또는 $N$은 $M$으로 나누어진다. $M$이 $N$과 다른 원소들을 포함하면 $M$은 $N$의 진약수가 된다. $N\eqzero{M}$, $M\eqzero{N}$이면 $M=N$이다. 최대공약과 최소공배수의 정의도 글자 그대로 유지된다. 특히 모듈 $M$이 유한개의 원소 $\alpha_1,\ldots,\alpha_e$를 포함하고 $M=(\alpha_1\ldots\alpha_e)$, 즉 모든 $\alpha\eqzero{M}$에 대하여
\[
 \alpha=c_1\alpha_1+\cdots+c_e\alpha_e+n_1\alpha_1+\cdots+n_e\alpha_e
\]
가 되며, 여기서 $c_i$가 $\Sigma$의 원소이고 $n_i$가 정수이면, $M$을 \emph{유한모듈}, $\alpha_1,\ldots,\alpha_e$를 \emph{모듈기저}라고 한다.

\emph{이하에서는 \S~1과 마찬가지로 다음 유한성 조건을 만족하는 환영역 $(\Sigma,T)$만을 바탕으로 삼는다. $(\Sigma,T)$의 모든 모듈은 유한모듈, 따라서 모듈기저를 가진다.}

그러면 정리 I의 유한사슬 정리는 그곳의 증명이 보여 주듯 이 환영역 $(\Sigma,T)$의 모듈에도 성립하고, 따라서 \S\S~2, 3의 모든 전제가 충족된다. 최단표현과 축약표현에 관한 정의 I과 보조정리 I을 직접 옮길 수 있다. 각 모듈을 유한개의 기약모듈의 최소공배수로 나타낼 수 있다는 정리 II도 마찬가지이며, 보조정리 II는 그러한 최단표현이 모두 동시에 축약되어 있음을 보여 준다. 모듈의 가약성을 그 보충부분의 성질로 나타내는 정리 III도 유지되고, 여기서 보조정리 III과 마침내 \emph{한 모듈을 기약모듈들의 최소공배수로 나타내는 서로 다른 두 최단표현에서 성분 수가 같다}는 정리 IV가 따른다. 또한 정리 IV는 축약표현에 관한 보조정리 I의 역인 보조정리 IV를 준다.

\emph{따름정리.} 같은 논증은 양쪽 환영역 $T$, 즉 우측이면서 좌측인 환영역에 대하여 모듈을 일관되게 \emph{양쪽모듈}, 곧 $\alpha$와 함께 $c\cdot\alpha$와 $\alpha\cdot c$를 모두 포함하고 $\alpha,\beta$와 함께 $\alpha-\beta$를 포함하는 모듈로 이해해도 이 모든 정리와 정의가 유지됨을 보여 준다.

이 모든 정리는 \emph{나눗셈}과 \emph{최소공배수}라는 개념에만 의존하지만, 뒤의 유일성 정리들은 본질적으로 \emph{곱}의 개념에 의존하므로 직접 옮길 수 없다. 예를 들어 $T$의 두 원소의 곱이 정의되어 있지 않으므로 으뜸아이디얼과 소아이디얼의 정의를 모듈로 옮길 수 없다. 형식적으로 가능한 비가환환으로의 이식도 의미를 잃는다. 여기서는 대응 소아이디얼의 존재를 증명할 수 없고\footnote{실제로 여기서는 $p_1^{\lambda_1}\eqzero{\ideal Q}$, $p_2^{\lambda_2}\eqzero{\ideal Q}$로부터 $(p_1-p_2)^{\lambda_1+\lambda_2}\eqzero{\ideal Q}$가 따르지 않으며, $(a\cdot b)^\lambda\eqzero{\ideal Q}$로부터 $a^\lambda\cdot b^\lambda\eqzero{\ideal Q}$도 따르지 않는다. 따라서 $\ideal P$가 아이디얼임도 증명할 수 없고, 소아이디얼의 성질을 가짐도 증명할 수 없다. 양쪽아이디얼의 경우에는 $\ideal P$가 아이디얼임은 증명할 수 있지만 소아이디얼임은 증명할 수 없다.} 기약아이디얼이 으뜸임을 보이는 증명도 실패하기 때문이다. 그런데 이 두 사실이 뒤의 유일성 정리들의 토대이다. 반면 비가환환이 단위원을 가지면 서로소 및 서로소-기약아이디얼을 정의할 수 있고, 정리 II의 증명에 사용한 논증은 각 아이디얼을 \emph{유한개의} 쌍마다 서로소인 서로소-기약아이디얼들의 최소공배수로 나타낼 수 있음을 보여 준다.\footnote{특수한 ``완전가약'' 아이디얼에 대해서는 여기서도 유일성 정리를 세울 수 있다. 뇌터--슈마이들러의 공동논문을 참조하라.}

끝으로 $(\Sigma,T)$가 유한성 조건을 만족하기 위한 충분조건을 하나 덧붙인다. \emph{$\Sigma$가 단위원을 가지고 유한성 조건을 만족하며, $T$ 자체가 $(\Sigma,T)$의 유한모듈이면 $(\Sigma,T)$의 모든 모듈은 유한모듈이다.}

실제로 $\Sigma$에 단위원이 존재하므로
\[
 \ideal M=(f_1,\ldots,f_e),\qquad
 f=\bar b_1f_1+\cdots+\bar b_e f_e+n_1f_1+\cdots+n_e f_e,
\]
여기서 $\bar b_i$는 $\Sigma$의 원소이고 $n_i$는 정수인 표현은
\[
 f=b_1f_1+\cdots+b_e f_e
\]
라는 표현도 준다. 여기서 $b_i$는 $\Sigma$의 원소이다. 실제로 $f_i=\varepsilon f_i$이므로 $n_if_i=n_i\varepsilon f_i$이고, $n_i\varepsilon=(\varepsilon+\cdots+\varepsilon)$는 $\Sigma$에 속하므로 $b_i=\bar b_i+n_i\varepsilon$도 $\Sigma$의 원소이다. 이제 $T$에 관한 가정은 $T$의 모든 원소 $\alpha$가
\[
 \alpha=\bar a_1\tau_1+\cdots+\bar a_k\tau_k
       +n_1\tau_1+\cdots+n_k\tau_k
\]
꼴의 표현, 따라서
\[
 \alpha=a_1\tau_1+\cdots+a_k\tau_k
\]
꼴의 표현을 가진다는 뜻이다. 둘째 표현은 $\varepsilon\tau_i=\tau_i$이므로 앞에서와 같이 첫째 표현에서 얻어진다.

이제 $\alpha$가 $(\Sigma,T)$의 한 모듈 $M$의 원소들을 두루 달리면, $\tau_k$의 계수 $a_k$들은 $\Sigma$의 한 아이디얼 $\ideal M_k$를 두루 달린다. 따라서 위의 결과에 따라 모든 $a_k\eqzero{\ideal M_k}$는
\[
 a_k=b_1a_k^{(1)}+\cdots+b_ea_k^{(e)}
\]
로 나타난다. $\alpha^{(i)}$가 $M$의 원소이고 그 $\tau_k$의 계수가 $a_k^{(i)}$라고 하자. 그러면 $\alpha-b_1\alpha^{(1)}-\cdots-b_e\alpha^{(e)}$는 $M$에 속하면서 $\tau_1,\ldots,\tau_{k-1}$에만 의존하는 원소이다. $\tau_k$를 더 이상 포함하지 않는 이 원소들 전체는 모듈 $M'$을 이루며, 여기에 같은 절차를 반복할 수 있다. 유한번 반복하면 주장이 증명된다.

\begin{center}
\S~10.\\[0.25em]
\textbf{다항식 환영역이라는 특수한 경우.}
\end{center}

1. 바탕 환영역 $\Sigma$가 임의의 복소수 계수를 갖는 $x_1,\ldots,x_n$의 모든 다항식으로 이루어진다고 하자. 이 환영역은 힐베르트의 모듈기저 정리(\emph{Math. Ann.} 36)에 따라 유한성 조건을 만족한다. 여기서는 \emph{우리 정리들과 소거이론 및 모듈론의 알려진 정리들 사이의 관계}를 다룬다.

이 관계는 잘 알려진 힐베르트 정리의 다음 특수한 경우로 얻어진다.\footnote{``완전 불변식계에 관하여'', \emph{Math. Ann.} 42 (1893), \S~3, 313쪽.}

$f$가 소아이디얼 $\ideal P$의 모든 다항식의 영점인 $x_1,\ldots,x_n$의 모든 유한 값계($\ideal P$의 영점)에서 0이 되면 $f$는 $\ideal P$로 나누어진다. 다시 말해 소아이디얼 $\ideal P$는 그 영점들에서 0이 되는 다항식 \emph{전체}로 이루어진다.\footnote{라스커[\emph{Math. Ann.} 60 (1905), 607쪽]가 보였듯 이 특수한 경우는 동차형식의 경우 직접 증명할 수도 있다. 거꾸로 여기서 동차 및 비동차 경우의 힐베르트 정리가 다시 따르며, 이를 다음과 같이 말할 수 있다. 아이디얼 $\ideal R$이 $\ideal M$의 모든 영점에서 0이 되면 $\ideal R$의 어떤 거듭제곱은 $\ideal M$으로 나누어진다. 그러나 이 정리와 위의 특수한 경우는 $x$의 값영역이 \emph{대수적으로 닫혀 있을 때}에만 성립한다. 따라서 우리 정리들만으로는 따를 수 없고 근의 존재를 이용해야 한다. 이를테면 영점이 $x_1=0,\ldots,x_n=0$뿐인 아이디얼은 어떤 차수 이상인 $x$들의 모든 멱곱을 포함한다는 대목에서 그러하다. 나머지 증명은 우리 정리들을 이용하면 라스커의 증명보다 조금 단순해진다. 라스커는 한 아이디얼을 극대 으뜸아이디얼들로 분해하는 일에도 힐베르트 정리를 사용해야 했기 때문이다.}

따라서 곱 $f\cdot g$가 $\ideal P$의 모든 영점에서 0이면 적어도 한 인자가 그 모든 영점에서 0이 된다. 그 영점들은 \emph{기약 대수도형}을 이룬다. 거꾸로 이것을 기약도형의 정의로 삼으면, 한 기약도형에서 0이 되는 모든 다항식은 소아이디얼을 이룬다. 따라서 소아이디얼과 기약도형은 일대일로 대응한다. 또한 $\ideal Q\eqzero{\ideal P}$, $\ideal P^e\eqzero{\ideal Q}$이므로 으뜸아이디얼의 영점들은 그에 대응 소아이디얼의 영점들과 일치한다.\footnote{으뜸아이디얼이 기약도형 하나를 가진다는 이 성질을 매콜리는 정의로 삼는다(서론 참조). 반면 라스커는 도형의 ``다양도''라는 개념만 정의에 넣고, 나머지는 \emph{추상적으로} 정의한다. $x_1=0,\ldots,x_n=0$에서만 0이 되는 으뜸아이디얼은 라스커에게 특수한 지위를 가진다.} 그러므로 한 아이디얼을 극대 으뜸아이디얼들의 최소공배수로 나타내면 그 아이디얼의 모든 영점을 기약도형들로 분해하게 되며, 라스커가 보였듯 그 역도 성립한다. \emph{따라서 대응 소아이디얼의 유일성을 증명하는 것은 여기서 대수도형을 기약도형들로 유일하게 분해할 수 있다는 소거이론의 기본정리에 대응한다.} 환영역 안의 다항식들을 기약다항식들의 곱으로 유일하게 나타낼 수 없고 따라서 소거이론도 존재하지 않는 특수한 다항식 환영역에서는, 우리의 정리를 그 소거이론 정리의 동치물로 볼 수 있다.

고립된 으뜸아이디얼들에 대응하는 기약도형들은 정확히 ``최소 레졸벤트''\footnote{예를 들어 J. König, \emph{Einleitung in die allgemeine Theorie der algebraischen Größen}(라이프치히, 토이프너, 1903), 235쪽을 참조하라.}에 나타나는 것들이다. 실제로 고립되지 않은 각 극대 으뜸아이디얼의 영점들은 적어도 하나의 고립아이디얼의 영점이기도 하다. 곧 그 고립아이디얼에 대응 소아이디얼은 고립되지 않은 아이디얼에 대응 소아이디얼로 나누어진다. 따라서 고립된 으뜸아이디얼의 유일성은 그 지수들에서 새로운 불변 중복도수를 준다. 으뜸아이디얼들을 기약아이디얼들로 분해하는 데 관한 유일성 정리들도 중복도의 관점에서 소거이론을 보완하는 것으로 볼 수 있다.

이 관찰들에 따라 여러 분해가 대수도형에 대하여 갖는 의미를 해석할 수 있다. 쌍마다 서로소인 아이디얼들에는 공통영점을 갖지 않는 도형들이 대응한다. 서로 상대소인 아이디얼들의 경우 한 아이디얼의 어느 기약도형도 다른 아이디얼의 공통영점도형이 아니다. 극대 으뜸아이디얼들은 모두 서로 다른 기약도형에서만 0이 된다. 기약아이디얼로 분해할 때에는 같은 기약도형이 여러 번 나타날 수도 있다.

일반 다항식 환영역 대신 모든 동차형식의 환영역을 바탕으로 삼아도 된다는 점을 덧붙인다. 거기서는 같은 차수의 형식들에 대해서만 덧셈이 정의되지만, 이 연산들 아래에서도 일반 정리들이 유지됨을 쉽게 확인할 수 있다.\footnote{그러나 다의성이 있는 경우 동차분해뿐 아니라 비동차분해도 존재할 수 있음은 다음 예가 보여 준다. $(x^3,xy,y^3)=[(x^3,y);(y^3,x)]=[(xy,x^3,y^3,x+y^3);(xy,x^3,y^3,y+x^2)]$.}

서로 다른 네 분해의 가장 간단한 예 하나는, 아래에 공식을 적을 것인데, 위의 논의에 따라 한 직선과 그 직선과 각각 꼬인 위치에 있으면서 서로 교차하는 두 직선으로 주어진다. 이 두 직선 가운데 하나는 두 직선의 교점과 다른 한 점을 더 높은 중복도로 포함한다. 서로소-기약아이디얼 분해에는 첫 직선과 그에 꼬인 위치에 있는 도형으로의 분해가 대응한다. 이 도형은 상대소-기약아이디얼 분해에서 두 직선으로 나뉜다. 극대 으뜸아이디얼 분해에는 더 높은 중복도의 점을 분리해 내는 일이 대응하고, 기약아이디얼 분해에는 이 점의 분해가 수반된다.

이 점을 원점으로, 그 점을 지나는 직선을 $y$축으로, 이 직선과 교차하는 직선을 $x$축에 평행하게, 그에 꼬인 위치의 직선을 $z$축에 평행하게 잡으면, 그러한 배치는 다음 \emph{기약아이디얼들}로 나타낼 수 있다.\footnote{이 가운데 처음 세 아이디얼은 소아이디얼이므로 기약이다. $\ideal B_4$는 약수 $(x^2,y,z)$와 $(x,y,z)$만 가지므로 기약이고, $\ideal B_5$는 그 모든 약수가 다항식 $xy$를 포함하므로 기약이다.}
\par\medskip
\noindent\makebox[\linewidth][c]{\(\displaystyle
 \ideal B_1=(x-1,y);\qquad
 \ideal B_2=(y-1,z);\qquad
 \ideal B_3=(x,z);\qquad
 \ideal B_4=(x^3,y,z);
\)}\par
\noindent\makebox[\linewidth][c]{\(\displaystyle
 \ideal B_5=(x^2,y^2,z).
\)}\par\medskip
\emph{대응 소아이디얼들}은 다음과 같다.
\par\smallskip
\noindent\makebox[\linewidth][c]{\(\displaystyle
 \ideal P_1=\ideal B_1;\qquad
 \ideal P_2=\ideal B_2;\qquad
 \ideal P_3=\ideal B_3;\qquad
 \ideal P_4=\ideal P_5=(x,y,z).
\)}\par\medskip
\emph{극대 으뜸아이디얼들}은 다음과 같다.
\par\smallskip
\noindent\makebox[\linewidth][c]{\(\displaystyle
 \ideal Q_1=\ideal B_1;\qquad
 \ideal Q_2=\ideal B_2;\qquad
 \ideal Q_3=\ideal B_3;\qquad
 \ideal Q_4=[\ideal B_4,\ideal B_5]=(x^3,y^2,x^2y,z).
\)}\par\medskip
\emph{상대소-기약아이디얼들}은 다음과 같다.
\par\smallskip
\noindent\makebox[\linewidth][c]{\(\displaystyle
 \ideal R_1=\ideal Q_1;\qquad
 \ideal R_2=\ideal Q_2;\qquad
 \ideal R_3=[\ideal Q_3,\ideal Q_4]=(x^3,x^2y,xy^2,z).
\)}\par\medskip
\emph{서로소-기약아이디얼들}은 다음과 같다.
\par\smallskip
\noindent\makebox[\linewidth][c]{\(\displaystyle
 \ideal S_1=\ideal R_1;\qquad
 \ideal S_2=[\ideal R_2,\ideal R_3]
 =((y-1)x^3,(y-1)x^2y,(y-1)xy^2,z).
\)}\par\medskip
이로부터 \emph{전체아이디얼}은
\par\smallskip
\noindent\(\displaystyle \ideal M=[\ideal S_1,\ideal S_2]=\ideal S_1\cdot\ideal S_2\)\par
\noindent\(\displaystyle
 =((x-1)(y-1)x^3,(y-1)x^2y,(y-1)xy^2,(x-1)z,y(y-1)x^3,yz),
\)\par\smallskip
\noindent\begin{tabular*}{\linewidth}{@{}l@{\extracolsep{\fill}}c@{}}
여기서 & \(\displaystyle 1=-(y-1)x^3+(y-1)(x^3-1)+y,\)
\end{tabular*}\par
\noindent\makebox[\linewidth][c]{\(\displaystyle
 (y-1)(x^3-1)+y\eqzero{\ideal S_1};\qquad
 -(y-1)x^3\eqzero{\ideal S_2}
\)}\par
이다. 이때 아이디얼 $\ideal B_1,\ideal B_2,\ideal B_3$은 고립되어 있으므로 기약아이디얼 및 극대 으뜸아이디얼 분해에서도 유일하게 정해진다. 아이디얼 $\ideal B_4,\ideal B_5$와 $\ideal Q_4$는 고립되어 있지 않다. 이들은 유일하게 정해지지 않으며, 예를 들어
\par\smallskip
\noindent\makebox[\linewidth][c]{\(\displaystyle
 \ideal D_4=(x^3,y+\lambda x^2,z);\qquad
 \ideal D_5=(x^2+\mu xy,y^2,z)
\)}\par
로 바꿀 수 있다. 마찬가지로 $\ideal Q_4$는
\par\smallskip
\noindent\makebox[\linewidth][c]{\(\displaystyle
 \overline{\ideal Q}_4=(x^3,x^2y,y^2+\lambda xy,z)
\)}\par
로 바꿀 수 있다.

2. 일반 다항식 환영역과 정수계수 다항식 환영역뿐 아니라, 다항식으로 이루어진 모든 \emph{유한 정역}도 힐베르트의 모듈기저 정리가 보여 주듯 유한성 조건을 만족한다.\footnote{거꾸로 한 다항식 환영역이 유한성 조건을 만족하고, 각 다항식이 $x$에 관하여 더 낮은 차수의 승수를 사용하는 표현을 적어도 하나 허용하면 그 환영역은 유한 정역이다.} 이때 계수들은 임의의 체에 속하게 할 수 있다. 여기서는 \emph{모든 짝다항식의 환영역}에 대한 예를 하나 더 제시한다. 이는 $x^2\cdot y^2=(xy)^2$ 때문에 환영역의 다항식들을 환영역의 기약다항식들의 곱으로 유일하게 나타낼 수 없는 가장 간단한 환영역이다. 앞의 예와 같은 배치를 이제 다음 \emph{기약아이디얼들}로 나타낸다.
\par\medskip
\noindent\makebox[\linewidth][c]{\(\displaystyle
 \ideal B_1=(x^2-1,xy,y^2,yz);\qquad
 \ideal B_2=(y^2-1,xz,yz,z^2);
\)}\par
\noindent\makebox[\linewidth][c]{\(\displaystyle
 \ideal B_3=(x^2,xy,xz,yz,z^2);\qquad
 \ideal B_4=(x^4,xy,y^2,xz,yz,z^2);
\)}\par
\noindent\makebox[\linewidth][c]{\(\displaystyle
 \ideal B_5=(x^2,y^2,xz,yz,z^2).
\)}\par\medskip
\emph{대응 소아이디얼들}은 다음과 같다.
\par\smallskip
\noindent\makebox[\linewidth][c]{\(\displaystyle
 \ideal P_1=\ideal B_1;\qquad
 \ideal P_2=\ideal B_2;\qquad
 \ideal P_3=\ideal B_3;\qquad
 \ideal P_4=\ideal P_5=(x^2,xy,y^2,xz,yz,z^2).
\)}\par\medskip
$\ideal P_1$이 소아이디얼이라는 사실은 환영역의 각 다항식이
\par\smallskip
\noindent\makebox[\linewidth][c]{\(\displaystyle
 f\equiv\varphi(z^2)+xz\,\psi(z^2)\;(\ideal P_1)
\)}\par\smallskip
꼴이 된다는 데서 따른다. 실제로
\par\smallskip
\noindent\makebox[\linewidth][c]{\(\displaystyle
 f_1\equiv\varphi_1(z^2)+xz\,\psi_1(z^2);\qquad
 f_2\equiv\varphi_2(z^2)+xz\,\psi_2(z^2)
\)}\par\smallskip
라고 하면 $f_1\cdot f_2\eqzero{\ideal P_1}$로부터 방정식들
\par\smallskip
\noindent\makebox[\linewidth][c]{\(\displaystyle
 \varphi_1\varphi_2+z^2\psi_1\psi_2=0;\qquad
 \varphi_2\psi_1+\varphi_1\psi_2=0
\)}\par
이 따르고, 여기서 $f_1\eqzero{\ideal P_1}$ 또는 $f_2\eqzero{\ideal P_1}$를 얻는다. 꼭 같은 방식으로 $\ideal P_2$가 소아이디얼임을 알 수 있다. 환영역의 각 다항식은 $\ideal P_3$을 법으로 $y^2$의 다항식과 합동이므로 $\ideal P_3$도 소아이디얼이다. $\ideal P_4$는 환영역의 모든 다항식으로 이루어진다. 따라서 $\ideal B_1,\ideal B_2,\ideal B_3$은 소아이디얼이므로 기약이다. 한편 $\ideal B_4$와 $\ideal B_5$는 각각 진약수 $\ideal P_4$ 하나만을 가지므로 반드시 기약이다.

기약아이디얼들로부터 다음 \emph{극대 으뜸아이디얼들}을 얻는다.
\par\medskip
\noindent\makebox[\linewidth][c]{\(\displaystyle
 \ideal Q_1=\ideal B_1;\qquad
 \ideal Q_2=\ideal B_2;\qquad
 \ideal Q_3=\ideal B_3;\qquad
 \ideal Q_4=[\ideal B_4,\ideal B_5]=(x^4,x^3y,y^2,xz,yz,z^2);
\)}\par\medskip
다음은 \emph{상대소-기약아이디얼들}이다.
\par\smallskip
\noindent\makebox[\linewidth][c]{\(\displaystyle
 \ideal R_1=\ideal Q_1;\qquad
 \ideal R_2=\ideal Q_2;\qquad
 \ideal R_3=[\ideal Q_3,\ideal Q_4]
 =(x^4,x^3y,x^2y^2,xy^3,xz,yz,z^2),
\)}\par\medskip
다음은 \emph{서로소-기약아이디얼들}이다.
\par\smallskip
\noindent\makebox[\linewidth][c]{\(\displaystyle
 \ideal S_1=\ideal R_1;\qquad
 \ideal S_2=[\ideal R_2,\ideal R_3]
\)}\par
\noindent\makebox[\linewidth][c]{\(\displaystyle
 =((y^2-1)x^4,(y^2-1)x^3y,(y^2-1)x^2y^2,(y^2-1)xy^3,xz,yz,z^2).
\)}\par\medskip
첫째 예에서와 마찬가지로
\par\smallskip
\noindent\makebox[\linewidth][c]{\(\displaystyle
 [\ideal B_4,\ideal B_5]=[\ideal D_4,\ideal D_5]
\)}\par
이고, 여기서
\noindent\(\displaystyle
 \ideal D_4=(x^4,xy+\lambda x^2,\ldots);\qquad
 \ideal D_5=(x^2+\mu xy,\ldots),\qquad \lambda\cdot\mu\ne1.
\)\par
또한
\noindent\(\displaystyle [\ideal Q_3,\ideal Q_4]=[\ideal Q_3,\overline{\ideal Q}_4],\)\quad
여기서 \(\displaystyle \overline{\ideal Q}_4=(x^4,x^3y,y^2+\lambda xy,\ldots).\)\par
나머지 기약아이디얼 또는 극대 으뜸아이디얼 $\ideal B_1,\ideal B_2,\ideal B_3$은 고립아이디얼이므로 유일하게 정해진다.

\begin{center}
\S~11.\\[0.25em]
\textbf{수론과 미분식 이론의 예.}
\end{center}

1. 바탕 환영역 $\Sigma$가 모든 \emph{짝수 유리정수}로 이루어진다고 하자. 그러면 $\Sigma$의 각 수 $2a$에 수 $a$를 대응시킴으로써 $\Sigma$를 모든 유리정수의 환영역과 일대일로 대응시킬 수 있다. 여기서 곧바로 $\Sigma$의 모든 아이디얼이 주아이디얼 $(2a)$임이 따른다. 아이디얼의 각 원소 $2c$에 대한 기저표현 $2c=n\cdot2a$에서 홀수 $n$은 유한합을 줄여 쓰는 기호일 뿐이다.

이 환영역의 \emph{소아이디얼들}은 $\ideal P_0=\ideal D=(2)$와 $\ideal P=(2p)$로 주어지며, 여기서 $p$는 홀수 소수이다. 따라서 각 소아이디얼은 $\ideal P_0$로 나누어지지만 다른 어느 소아이디얼로도 나누어지지 않는다. \emph{으뜸아이디얼들}은
\[
 \ideal Q_{e_0}=(2\cdot2^{e_0})
 \qquad\text{및}\qquad
 \ideal Q_e=(2p^e)
\]
로 주어진다. 이들은 동시에 \emph{기약아이디얼들}이다. 소아이디얼들에 관하여 방금 말한 사실에 따라 서로 다른 홀수 소수에 속하는 어느 두 아이디얼도 서로 상대소이지만, 어느 $\ideal Q$도 어떤 $\ideal Q_{e_0}$에 대해서는 상대소가 아니다.\footnote{실제로 홀수 $p_1\ne p_2$에 대하여 $2b\cdot2p_1^{e_1}\eqzero{2p_2^{e_2}}$이면 언제나 $2b\eqzero{2p_2^{e_2}}$이다. 그러나 $2b\cdot2p^e\eqzero{2\cdot2^{e_0}}$에서는 $2b\equiv2\cdot2^{e_0-1}\,(2\cdot2^{e_0})$만이 따른다.} 따라서 $\ideal Q_{e_0}$들이 유일한 비고립 으뜸아이디얼들이다. $a$를 소수 거듭제곱들로 유일하게 분해하는 것에는 아이디얼 $(2a)$을 극대 으뜸아이디얼이면서 동시에 기약아이디얼인 것들로 나타내는 \emph{유일한} 표현
\[
 (2a)=[(2\cdot2^{e_0}),(2p_1)^{e_1},\ldots,(2p_\alpha)^{e_\alpha}]
\]
이 대응한다. 따라서 다항식 환영역의 예들과 달리 여기서는 \emph{비고립} 극대 으뜸아이디얼들도 \emph{유일하게} 정해진다. $e_0=0$이면 이 표현은 동시에 서로 상대소인 아이디얼들로 된 표현이고, $e_0>0$이면 그 아이디얼은 상대소-기약이다.

그러므로 \emph{모든 정수}의 환영역에서는 서로 다른 네 분해가 일치하지만, 여기서는 한편으로 극대 으뜸아이디얼 분해와 기약아이디얼 분해만이 일치한다. 다른 한편 $e_0>0$일 때에는 서로소-기약 분해(환영역이 단위원을 갖지 않으므로 모든 아이디얼이 서로소-기약이다)와 상대소-기약 분해가 일치하지만, $e_0=0$일 때에는 서로소-기약 분해와 상대소-기약 분해가 서로 달라진다. 동시에 여기서는 한 소아이디얼이 다른 소아이디얼로 나누어지면서도 그와 같지 않을 수 있다는 예, 더 일반적으로 \emph{나눗셈으로부터 곱표현이 따르지 않는다는} 예를 이미 얻는다. 후자의 사실은 환영역에 단위원이 없다는 데서 생기며, 모든 아이디얼이 주아이디얼임에도 환영역의 수들을 환영역의 기약수들의 곱으로 유일하게 나타낼 수 없는 이유이기도 하다. 따라서 여기서는 최소공배수를 도입하는 일이 필요하다. 모든 짝수 대신 \emph{고정된 소수 또는 소수 거듭제곱으로 나누어지는 모든 수}를 바탕으로 삼아도 사정은 정확히 \emph{똑같다}는 점을 덧붙인다.

반면 $\Sigma$가 \emph{합성수} $g=p_1^{\sigma_1}\cdots p_\nu^{\sigma_\nu}$로 나누어지는 모든 수로 이루어지면 \emph{기약아이디얼과 으뜸아이디얼도 서로 달라진다.} 다시 모든 아이디얼은 주아이디얼 $(g\cdot a)$이다. \emph{소아이디얼들}은 다시 $\ideal P_0=\ideal D=(g)$와 $\ideal P=(g\cdot p)$로 주어지는데, 여기서 $p$는 $g$를 나누는 소수들과 다른 소수이다. 반면 \emph{기약아이디얼들}은 $\ideal B_{\lambda_i}=(g\cdot p_i^{\lambda_i})$와 $\ideal B_e=(g\cdot p^e)$로 주어진다. \emph{으뜸아이디얼들}은 $\ideal Q_e=\ideal B_e$와, 기약아이디얼들과는 다른
\[
 \ideal Q_{\lambda_1\ldots\lambda_\nu}
 =(g\cdot p_1^{\lambda_1}\cdots p_\nu^{\lambda_\nu})
\]
로 주어진다. 여기서 $\ideal B_{\lambda_i}$들과 $\ideal Q_{\lambda_1\ldots\lambda_\nu}$는 모두 같은 대응 소아이디얼 $\ideal P_0=(g)$를 가진다. 여기서도 기약아이디얼로의 분해는 유일하고, 따라서 극대 으뜸아이디얼로의 분해도 유일하다. 그러므로 여기서도 비고립아이디얼들까지 유일하게 정해진다.

2. \emph{비가환 환영역}의 한 예는 뇌터--슈마이들러의 논문에서 다룬 비가환 다항식 환영역의 아이디얼론이다. 특히 그 논문은 ``완전가약'' 아이디얼, 즉 분해의 성분들이 쌍마다 서로소이고 진약수를 갖지 않는 아이디얼을 다룬다. 따라서 성분들은 더더욱 기약이다. 그러므로 거기서 증명한 동형성에 더하여 \S~9로부터 서로 다른 두 분해에서 \emph{성분 수가 같다}는 결과를 얻는다. 이로써 그 논문의 특수한 경우로 나오는 편미분 또는 상미분 선형미분식들의 연립계 분해에 관한 결과를 얻는데, 이 결과는 상미분 선형미분식 하나라는 잘 알려진 경우에도 지적된 적이 없는 듯하다.

동시에 고정된 아이디얼 $\ideal M$의 모든 잉여류로 이루어진 체계 $T$는 비가환 다항식 환영역 $\Sigma$와 함께 이중환영역 $(\Sigma,T)$를 이루며, $T$는 $\Sigma$에 관한 모듈 성질을 가진다. 실제로 두 잉여류의 차는 다시 잉여류이고, 한 잉여류와 임의의 다항식의 곱도 잉여류이지만 두 잉여류의 곱은 존재하지 않는다(앞의 논문 \S~3). 그곳에서 ``부분군''이라고 부른 잉여류의 체계들은 따라서 바탕 환영역 $\Sigma$가 비가환인 이중환영역 $(\Sigma,T)$의 모듈 예를 이룬다.

\begin{center}
\S~12.\\[0.25em]
\textbf{초등인자론의 예.}
\end{center}

여기서는 일반 전개에서 도출되는 초등인자론의 관점을 제시하지만, 초등인자론 자체는 \emph{알려져 있다}고 가정한다.

$\Sigma$를 $n^2$개의 원소를 가진 모든 정수행렬의 환영역이라 하고, 덧셈과 곱셈은 행렬에 통상적인 뜻으로 정의하자. 그러면 $\Sigma$는 \emph{비가환 환영역}이다. 따라서 아이디얼들은 일반적으로 한쪽이고, 양쪽아이디얼들은 그 특수한 경우로만 포함된다.\footnote{이 환영역의 아이디얼론은 뒤 파스키에의 다음 논문들에서 다룬다. ``테타리온의 수론'', 취리히 학위논문, \emph{Vierteljahrsschr. d. Naturf. Ges. Zürich} 51 (1906); ``테타리온아이디얼의 이론에 관하여'', 같은 곳 52 (1907). 둘째 논문의 내용은 모든 아이디얼이 주아이디얼임을 증명하는 것이다.}

먼저 \emph{모든 아이디얼이 주아이디얼임}을 보이겠다. 이를 위하여 우측아이디얼의 경우 각 행렬 $A=(a_{ik})$에 정수계수 선형형식들의 모듈
\[
 A=(a_{11}\xi_1+\cdots+a_{1n}\xi_n,\ldots,
 a_{n1}\xi_1+\cdots+a_{nn}\xi_n)
\]
을 대응시킨다. 거꾸로 이 모듈에는 $A$의 기저를 주는 모든 행렬, 곧 $A$뿐 아니라 $U$가 유니모듈러인 $UA$도 대응한다. 더 일반적으로 곱 $PA$에는 $A$의 배수인 모듈 $B$가 대응한다. $A$의 한 선형형식은 0이 아닌 행을 하나만 가진 $P$들로 주어진다.

이제 $A_1,A_2,\ldots,A_\nu,\ldots$가 아이디얼 $\ideal M$의 모든 원소이고, 같은 기호 $A_1,A_2,\ldots,A_\nu,\ldots$가 그에 대응하는 모듈들이라고 하자. 이 모듈들의 최대공약을 $A$라 하고, 이 모듈 $A$에 대응하는 가장 일반적인 행렬을 $UA$라 하자. 최대공약의 정의에 따라 $A$의 각 선형형식에는 행렬 $P_1A_{i_1}+\cdots+P_\sigma A_{i_\sigma}$가 대응하며, 위에서와 같이 각 $P$는 0이 아닌 행을 하나만 가진다. 따라서 $A$의 한 기저에 대응하는 행렬 $A$도 이제 일반적인 $P$를 사용한 그러한 표현을 가지며, 그 결과 $\ideal M$의 원소가 된다. 더 나아가 각 모듈 $A_i$는 $A$로 나누어지므로 각 행렬 $A_i$도 $A$로 나누어진다. \emph{따라서 $A$, 더 일반적으로 $UA$가 $\ideal M$의 기저를 이룬다.} 좌측아이디얼의 경우에는 그에 맞추어 각 행렬의 열들을 한 모듈의 기저로 보아야 한다. 모든 아이디얼은 주아이디얼이고, $A$와 함께 $AV$도 그 기저가 되며, 여기서 $V$는 임의의 유니모듈러 행렬이다.

이제 초등인자론과 연결하기 위하여 이하에서는 \emph{양쪽아이디얼}을 바탕으로 삼는다. 따라서 특히 $A$와 함께 $PAQ$도 아이디얼에 속한다. 그러면 위의 결과에 따라 그러한 아이디얼의 가장 일반적인 기저는 $UAV$로 주어지며, 여기서 $U,V$는 유니모듈러이다. 따라서 기저원소들은 \emph{서로 동치인 행렬들의 한 류}를 모두 이룬다.\footnote{한쪽아이디얼의 기저원소들에는 각각 우측류 또는 좌측류가 대응한다.} 아이디얼과 류 사이에는 일대일 관계가 있고, 따라서 아이디얼과 그 류의 초등인자계 $(a_1\mid a_2\mid\cdots\mid a_n)$ 사이에도 일대일 관계가 있다. 여기서 $a_i$들은 잘 알려진 대로 음이 아닌 정수이며 각 수는 그 다음 수를 나눈다. 초등인자들이 정하는 표준형에 나타나는 그 류의 행렬을 따라서 아이디얼의 특별한 기저로 볼 수 있다. 아이디얼 또는 류의 나눗셈 관계에서 초등인자들의 나눗셈 관계가 따르고, 그 역도 성립한다.

그런데 뒤 파스키에는 앞의 논문 \S~11에서 모든 양쪽아이디얼의 랭크가 $n$이고 모든 초등인자가 일치함을 보였다. 따라서 일반적인 초등인자를 다루려면 아이디얼이 아니라 양쪽류 자체에서 직접 출발해야 한다(이 류들도 큰 독일문자로 나타낸다). 이때 한 류 $\mathfrak A=UAV$는 $A=PBQ$가 될 때 다른 류 $\mathfrak B$로 나누어진다.

일반적으로 다음이 성립한다. \emph{두 류의 최소공배수(최대공약)는 대응하는 초등인자계들의 최소공배수(최대공약)를 취하여 얻는다.}\footnote{E. 슈타이니츠는 ``모듈론에 관하여'', \emph{Math. Ann.} 52 (1899), 1쪽에서 초등인자계들의 최소공배수와 최대공약으로 류의 최소공배수와 최대공약을 정의한다. 이와 독립적으로 류의 최소공배수는 H. 브란트, ``한 기본형식에 관한 이차 이항형식의 합성'', \emph{J. f. M.} 150 (1919), 1쪽에 ``합동합성''이라는 이름으로 나타난다.} 실제로 $(a_1\mid a_2\mid\cdots\mid a_n)$, $(b_1\mid b_2\mid\cdots\mid b_n)$, $(c_1\mid c_2\mid\cdots\mid c_n)$가 각각 $\mathfrak A,\mathfrak B,\mathfrak C^*$의 초등인자계이고 $c_i=[a_i,b_i]$라고 하자. 또한 $\mathfrak C=[\mathfrak A,\mathfrak B]$라 하자. 그러면 $\mathfrak C^*$는 $\mathfrak A$와 $\mathfrak B$로, 따라서 $\mathfrak C$로 나누어진다. 거꾸로 $\mathfrak C$의 초등인자들은 $\mathfrak C^*$의 초등인자들로 나누어지므로 $\mathfrak C$는 $\mathfrak C^*$로 나누어진다. 따라서 $\mathfrak C=\mathfrak C^*$이다. 최대공약의 증명도 그에 맞추어 얻어진다.

초등인자 $a_i$를 소수 거듭제곱들의 최소공배수로 유일하게 나타내는 것에는 $\mathfrak A$를 류 $\mathfrak Q$들의 최소공배수로 나타내는 표현이 대응한다. 이 $\mathfrak Q$들의 초등인자는 한 소수의 거듭제곱들로 주어지며, 기호로는
\[
 \mathfrak Q\sim(p^{r_1}\mid p^{r_2}\mid\cdots\mid p^{r_e}\mid0\mid\cdots\mid0);
 \qquad r_1\leqq r_2\leqq\cdots\leqq r_e
\]
이다. 특히 랭크 $\varrho$가 $n$과 같으면 이는 서로소인 서로소-기약류들로의 분해이며, 비가환 환영역임에도 이 분해는 유일하다.

류 $\mathfrak Q$는 다음 초등인자계들에 대응하는 류들로 더 분해할 수 있다.
\[
\begin{gathered}
 \mathfrak B_1\sim(p^{r_1}\mid\cdots\mid p^{r_1});\quad
 \mathfrak B_2\sim(1\mid p^{r_2}\mid\cdots\mid p^{r_2});\quad \ldots;\\
 \mathfrak B_e\sim(1\mid\cdots\mid1\mid p^{r_e}\mid\cdots\mid p^{r_e});\quad
 \mathfrak B_{e+1}\sim(1\mid\cdots\mid1\mid0\mid\cdots\mid0).
\end{gathered}
\]
여기서 $\mathfrak B_\nu$에는 매번 수 $1$이 $(\nu-1)$번 나타난다. 이를테면 $r_1=r_2=\cdots=r_\mu=0$이면 $\mathfrak B_1,\ldots,\mathfrak B_\mu$는 단위류와 같아지므로 분해에서 생략해야 한다. $\varrho=n$이면 $\mathfrak B_{e+1}$에도 같은 말이 성립한다. 더 나아가 이를테면 $r_\nu=r_{\nu+1}=\cdots=r_{\nu+\lambda}$이면 $\mathfrak B_{\nu+1},\ldots,\mathfrak B_{\nu+\lambda}$는 $\mathfrak B_\nu$의 진약수가 되므로 역시 제외해야 한다. 남은 류들을 $\mathfrak B_{i_1},\ldots,\mathfrak B_{i_k}$라 하면 이들은 이제 최단표현을 주고,
\[
 \mathfrak Q=[\mathfrak B_{i_1}\ldots\mathfrak B_{i_k}]
\]
는 \emph{$\mathfrak Q$를 기약류들로 나타내는 유일한 분해}이다.

실제로
\[
 \mathfrak B_\nu\sim(1\mid\cdots\mid1\mid p^{r_\nu}\mid\cdots\mid p^{r_\nu})
\]
를
\[
 \mathfrak C\sim(1\mid\cdots\mid1\mid p^{s_1}\mid\cdots\mid p^{s_\lambda})
 \quad\text{및}\quad
 \mathfrak D\sim(1\mid\cdots\mid1\mid p^{t_1}\mid\cdots\mid p^{t_\mu})
\]
의 최소공배수로 나타낼 수 있다고 하자. 그러면 $\mathfrak C$와 $\mathfrak D$의 초등인자계에서 처음 $(\nu-1)$자리에는 수 $1$이 놓여야 한다. $\nu$번째 자리에서는 지수 $s_1$ 또는 $t_1$, 이를테면 $s_1$이 $r_\nu$와 같아야 한다. 그런데
\[
 r_\nu=s_1\leqq s_2\leqq\cdots\leqq s_\lambda\leqq r_\nu
\]
이므로 $\mathfrak C=\mathfrak B_\nu$이다. 따라서 $\mathfrak B_\nu$는 \emph{기약}이다.\footnote{반면 $\mathfrak B$들은 한쪽류들로는 더 분해될 수 있다. 여기서는 초등인자와 류 사이의 관계가 더 이상 일대일이 아니고, 따라서 기약 한쪽류들로의 분해도 유일하지 않다. H. 브란트가 알려 준 다음 우측류 분해의 예가 이를 보여 준다. 여기서 류는 한 기저행렬 또는 그에 대응하는 모듈로 나타낸다.
\[
\begin{gathered}
 \mathfrak B=[\mathfrak C_1,\mathfrak C_2]=[\mathfrak D_1,\mathfrak D_2],\quad
 \mathfrak B\sim\begin{pmatrix}p&0\\0&p\end{pmatrix},\quad
 \mathfrak C_1\sim\begin{pmatrix}1&0\\0&p\end{pmatrix},\quad
 \mathfrak C_2\sim\begin{pmatrix}p&0\\0&1\end{pmatrix},\\
 \mathfrak D_1\sim\begin{pmatrix}1&0\\0&p\end{pmatrix},
 \quad
 \mathfrak D_2\sim\begin{pmatrix}p&0\\(p-1)&1\end{pmatrix}.
\end{gathered}
\]
실제로 모듈 $(\xi,p\eta)$, $(p\xi,\eta)$, $(p\xi,(p-1)\xi+\eta)$는 각각 진약수로 모듈 $(\xi,\eta)$만을 가지므로 기약이며, 더 나아가 서로 다르다.} $p$를 $0$으로 바꾸면 $\mathfrak B_{e+1}$에도 같은 사실이 성립한다. 그런데 이 기약류들은 각각 최소공배수를 만들 때 하나의 특정한 지수와 그 지수가 $\mathfrak Q$ 또는 $\mathfrak A$의 초등인자계에서 처음 나타나는 자리를 주며, $\mathfrak B_{e+1}$은 랭크를 준다. 이 수들은 $\mathfrak Q$ 또는 $\mathfrak A$의 초등인자계에 의해 유일하게 정해지고 초등인자와 류 사이의 관계는 일대일이므로, $\mathfrak Q$의 기약류 분해, 마찬가지로 임의의 류의 기약류 분해는 \emph{유일하다}.

정리하면 다음과 같다. \emph{원소 수가 제한된 정수행렬들의 모든 양쪽류 $\mathfrak A$는 유한개의 기약 양쪽류들의 최소공배수로 유일하게 나타낼 수 있다. 각 기약류는 $\mathfrak A$의 초등인자계의 고정된 소인수 하나, 그에 딸린 지수 하나, 그리고 그 지수가 처음 나타나는 자리를 나타낸다. 인수 $0$에 대응하는 기약류는 $\mathfrak A$의 랭크를 나타낸다.}

\begin{flushleft}
에를랑겐, 1920년 10월.
\end{flushleft}
\begin{center}
(1920년 10월 16일 접수.)
\end{center}

\clearpage
\setcounter{footnote}{0}
